\documentclass[oneside]{book}

\title{Appunti di Wireline and Wireless Networks}
\author{Laura Straccini}
\date{\today}

% --- INIZIO richiamo di pacchetti di utilità.
\usepackage[plainpages=false]{hyperref}	% generazione di collegamenti ipertestuali su indice e riferimenti
\usepackage[all]{hypcap} % per far si che i link ipertestuali alle immagini puntino all'inizio delle stesse e non alla didascalia sottostante
\usepackage{amsthm}	% per definizioni e teoremi
\usepackage{amsmath}	% per ``cases'' environment
% --- FINE riachiamo di pacchetti di utilità

\usepackage{xcolor}
\usepackage{listings}

\usepackage{float}
\usepackage{graphicx}

% Definizione colori stile "Vibrant Neon" (Sfondo schiarito e più colori)
\definecolor{darkslateblue}{RGB}{40, 40, 54}    % Sfondo grigio meno cupo
\definecolor{darkturquoise}{RGB}{0, 206, 209}     % Keyword principali
\definecolor{deeppink}{RGB}{255, 20, 147}  % Stringhe
\definecolor{gainsboro}{RGB}{220, 220, 220}    % Commenti
\definecolor{yellowneon}{RGB}{255, 255, 0}   % Tipi di dati (int, float, ecc.)
\definecolor{darkviolet}{RGB}{147, 50, 186}   % Numeri e Costanti
\definecolor{softwhite}{RGB}{240, 240, 240}  % Testo base (variabili)

\lstset{
    backgroundcolor=\color{darkslateblue},   
    commentstyle=\color{gainsboro}\itshape,
    keywordstyle=\color{darkturquoise}\bfseries,
    % Aggiungiamo enfasi per i tipi di dati comuni in C++
    emph={int,char,double,float,unsigned,void,bool},
    emphstyle=\color{yellowneon},
    numberstyle=\tiny\color{darkviolet},
    stringstyle=\color{deeppink},
    basicstyle=\ttfamily\scriptsize\color{softwhite}, % "scriptsize" lo rende più piccolo
    breaklines=true,                 
    captionpos=b,                    
    keepspaces=true,                 
    numbers=left,                    
    numbersep=8pt,                  
    showspaces=false,                
    showstringspaces=false,
    language=C++,
    frame=leftline,                 
    rulecolor=\color{deeppink},  % Cambiato in magenta per staccare dallo sfondo
    xleftmargin=12pt,
    % Per evidenziare i numeri nel codice
    literate={0}{{\textcolor{darkviolet}{0}}}{1}
             {1}{{\textcolor{darkviolet}{1}}}{1}
             {2}{{\textcolor{darkviolet}{2}}}{1}
             {3}{{\textcolor{darkviolet}{3}}}{1}
             {4}{{\textcolor{darkviolet}{4}}}{1}
             {5}{{\textcolor{darkviolet}{5}}}{1}
             {6}{{\textcolor{darkviolet}{6}}}{1}
             {7}{{\textcolor{darkviolet}{7}}}{1}
             {8}{{\textcolor{darkviolet}{8}}}{1}
             {9}{{\textcolor{darkviolet}{9}}}{1}
}


\begin{document}

\frontmatter
\maketitle % Sostituisce \generaFrontespizio
\tableofcontents % Sostituisce \generaIndice

\mainmatter

% Sostituiamo il comando \capitolo con i comandi standard \chapter e \input (o \include)
\chapter{Sistemi di Comunicazione Qualsiasi}
\section{Introduzione}
La prima classificazione dei segnali distingue i segnali tra analogici e digitali. 

I segnali \textbf{analogici} sono continui nel tempo (prevediamo l'andamento tra -inf e +inf) 
e possono assumere un numero infinito di valori (un esempio di segnale analofico è la voce
modifca l'aria, le particelle che compongono l'aria e si propaga nello spazio; un altro esempio 
i segnali radio e elettrocardiogramma), vengono trasemssi e sono riproducibili, il segnale 
analogico per eccellenza è il vinile, il solco riproduce perfettamente le frequenze.
Il loro problema è che sono sensibili al rumore, se c'è un disturbo, il segnale viene
alterato e non è più possibile recuperare il segnale originale, è necessario quindi un 
segnale più robusto (segnale digitale che posso ricostruire anche se c'è un disturbo, 
basta che il disturbo non sia troppo grande).
Quando io passo passo da un segnale analogico a uno digitale ho fatto il
\textbf{campionamento}, estraggo i campioni a istanti di tempo.
Quasi tutti i segnali che trasmettiamo sono digitali, ma i segnali analogici campioanti
ancora non sono digitali (sequenza di 0 e 1).
Il campionamento mi dice che è vero che perdi informazione, ma Shannon ci dice che è 
possibile ricostruire il segnale originale (quello che sta in mezzo e viene perso).

I segnali \textbf{digitali} sono discreti nel tempo e possono assumere solo un numero 
finito di valori.

\section{Concetti base sui segnali}

Per spigare come funziona il teorema di Nyquist-Shannon, è necessario introdurre alcuni
concetti base sui segnali.
\subsection{Segnali sinusoidali}

\textbf{Periodo}: Intervallo di tempo dopo cui il segnale si ripete, si misura in secondi.

\textbf{Frequenza}: Numero di periodi al secondo, si misura in Hertz (Hz), è l'inverso del periodo.

\textbf{Ampiezza}: Valore massimo del segnale, si misura in Volt (V).

\textbf{Fase}: Punto di partenza del segnale, si misura in gradi o radianti.

\subsection{Serie di Fourier}
Se io ho un segnale periodico, comunque sia fatto (esempio onda quadra periodica), 
si dimostra che tutti i segnali periodici con periodo T possono essere rappresentati 
come somma di segnali sinusoidali (seno e coseno) con frequenze multiple di 1/T, 
cioè con frequenze discrete.\\
Questa rappresentazione è chiamata \textbf{serie di Fourier} e ci dice che qualsiasi 
segnale periodico può essere scomposto in una serie di segnali sinusoidali, ognuno 
con una frequenza, ampiezza e fase specifica.

Di solito non trasmetto segnali periodici, ma serve per introdurre trasformata di fourier.

Lo \textbf{spettro di ampiezza} e lo \textbf{spettro di fase} di un segnale 
sono rispettivamente il modulo e la fase dei coeffienti di Fourier.\\
Questa rappresentazione nel dominio della frequenza è complementare alla rappresentazione
del dominio del tempo.\\
Se conosco lo spettro di ampiezza e di fase conosco il segnale.\\
Dimensionamento del canale: capire quanta banda mi serve per trasmettere un segnale.

Effetto Gibbs: quando approssimo un segnale con una serie di Fourier, ottengo 
un'oscillazione vicino ai punti di discontinuità, che non scompare aumentando 
il numero di termini della serie, ma si riduce a una zona sempre più piccola 
intorno alla discontinuità.

\subsection{Trasformata di Fourier}
La \textbf{trasformata di Fourier} estende il concetto della Serie di Fourier ai segnali
non periodici.
 In questo corso la trasformata di fourier si chiama spettro del segnale.

Propietà della trasformata: ...

Nessun sistema di trasmissione è perfetto, la priam cosa che fa è attenuare il segnale. 
Se tutte le componenti di Fourier sono attenuate allo stesso modo, il segnale è 
attenuato ma non distorto, se invece alcune componenti sono più attenuate di altre, 
il segnale è distorto.

Tutti i canali di trasmissione hanno una banda ed è caratteristico di ogni sistema.
Il teorema di Shannon lega il bitrate massimo alla banda.

Nel wifi si usa un canale con una banda piccola.

\subsection{Conversione Analogico Digitale}
Il sistema deve prima campionare quindi trasformare il segnale analogico in uno tempo discreto
estrando dei campioni, Nyquist dice la freqeunza minima.

Per trasformare il numero in un dato digitale applico la quantizzazione.

Teorema di Nyquist ci dice prendi il segnale, vedi le frequenze, la frequenza di campionamento
deve essere almeno il doppio della frequenza massima del segnale. Sotto questa soglia si genera
l'aliasing, una distorsione che rende impossibile ricostruire il segnale originale.

L'errore di quantizzazione è la differenza tra il segnale analogico originale e il 
segnale digitale quantizzato, è un errore che si introduce quando si approssima un 
segnale continuo con un numero finito di livelli discreti.

Risoluzione in bit: con n bit posso rappresentare $2^n$ livelli di quantizzazione, 
più bit ho, più livelli di quantizzazione posso rappresentare, minore è l'errore 
di quantizzazione, ma maggiore è la quantità di dati da trasmettere.

\subsection{Teorema del campionamento}
Quando campione devo campionare con una frequenza maggiore o uguale la frequezna di Nyquist
$f_c=2W$.

\subsection{Frequenze di Campionamento}
Per segnali audio, la frequenza di campionamento standard è di 44.1 kHz, 
che è più del doppio della frequenza massima udibile (20 kHz). 
Per segnali video, la frequenza di campionamento dipende dalla risoluzione 
e dalla frequenza dei fotogrammi, ma è generalmente molto più alta.

\subsection{Quantizzazione}
Una volta campionato, devo decidere quanti intervalli di quantizzazione usare, 
più intervalli uso, più preciso è il segnale digitale, ma più dati devo trasmettere.

L'attuale standard per il segnale audio è di 16 bit (cioè 65536 livelli).

Di quanto spazio (byte) abbiamo bisogno per memorizzare un'ora di musica?
Devo trasmettere ogni secondo 44100 campioni, ogni campione è rappresentato da 16 bit 
(2 byte) al secondo, i canali audio sono due (stereo), quindi 
$$44100 campioni/s * 2 byte/campione * 2 canali = 176400 byte/s$$
Byte=635MB da memorizzare

Il bitrate per trasmettere il segnale audio correttamente?
$$44100 campioni/s * 16 bit/campione * 2 canali = 1411200 bit/s = 1411.2 kbps$$

I byte vengono usati per memorizzare i dati, i bit al secondo vengono usati per indicare
la velocità di trasmissione dei dati.

\section{Companding Non Lineare}
Il companding è una tecnica di quantizzazione non lineare che migliora il rapporto 
segnale-rumore (SNR) per segnali di bassa ampiezza.

\section{PCM Pulse Code Modulation}
Si prende il segnale registrato dal microfono di un telefono e si dice che la banda è 4 kHz,
non è vero la banda della voce è arriva pure a 100 kHz, ma per risparmiare banda si taglia a 4 kHz,
quindi 4 kHz è la banda della voce, la moltiplico per 2 (teorema di Nyquist) e ottengo 8 kHz, 
quindi campiono a 8 kHz (8000 campioni/s), quindi un campione ogni 125 microsecondi, ogni campione 
lo quantizzo e lo codifico in 8 bit (256 livelli di quantizzazione), quindi il bitrate è di 64 kbps 
$(8000 campioni/s * 8 bit/campione)$. 

Devo filtrare il segnale prima di campionarlo, altrimenti rischio di avere aliasing, 
quindi uso un filtro a 3.4 kHz e poi campiono, estraggo un campione, questo campione viene
inviato in un quantizzatore che trasforma il segnale in un numero digitale a 8 bit.

Questo è il processo di \textbf{PCM (Pulse Code Modulation)}, è un processo di conversione 
analogico-digitale, vale per qualsiasi segnale, non solo per la voce, ma anche per il video, 
l'audio, ecc.\\

In ottica si chiama on off keying, è un processo di modulazione digitale,
dove il segnale digitale viene modulato su un segnale portante, in questo caso la luce, 
accendendo e spegnendo la luce per rappresentare i bit 1 e 0.\\

Trasformo una sequenza di bit in un segnale che posso trasmettere, si chiama \textbf{modulazione}.

Devo permettere l'accesso multiplo, al canale di comunicazione di accedere a più utenti, 
ci sono vari dominii dipende da cosa devo fare. 

Utente 1 deve trasmettere 8 k campioni /s, quindi un campione ogni 125 microsecondi;\\
Nel frattempo (nei 125 microsendi) trasmetto altri utenti, in particolare ci trasmetto 32 utenti.
Time division multiplexing (TDM): divido il tempo in slot, ogni slot è di 125 microsecondi,
e ogni utente trasmette in uno slot, quindi ogni utente trasmette ogni 4 ms $(32 slot * 125 microsecondi/slot)$,
quindi ogni utente trasmette 250 campioni/s, quindi un bitrate di 2 kbps $(250 campioni/s * 8 bit/campione)$.

\subsection{Codifica di Sorgente}
Cerca di \textbf{ridurre il numero di bit necessari per rappresentare un segnale}, sfruttando la
ridondanza del segnale, ad esempio usando la codifica di Huffman, che assegna codici più 
brevi ai simboli più frequenti.\\

Ci sono due tipi: quelli con la perdita di informazione (lossy) e quelli senza perdita 
di informazione (lossless).

Primo teorema di Shannon: l'entropia è 
$$H_s=E\{I_k\}=\sum_{k=1}^{L} p_k \log_2 \frac{1}{p_k}= \sum_{k=1}^{L} p_k I_k$$

\subsection{Codifica di Canale}
Aggiunge bit, ovvero aggiunge ridondanza, per permettere la rilevazione e la 
correzione degli errori, per proteggerloi dagli errori introdotti dal rumore o dalle 
interferenze durante la trasmissione. Il suo scopo è garantire l'integrità del messaggio.\\

Fa riferimento al secondo teorema di Shannon, che stabilisce un limite superiore alla 
capacità di un canale di comunicazione, indicando la massima velocità alla quale è 
possibile trasmettere informazioni con una probabilità di errore arbitrariamente bassa, 
dato un certo livello di rumore.\\

$$C = B \log_2 (1 + \frac{S}{N})$$
dove:
- $C$ è la capacità del canale in bit al secondo (bps)
- $B$ è la banda del canale in hertz (Hz)
- $S$ è la potenza del segnale in watt (W)
- $N$ è la potenza del rumore in watt (W)
Il secondo teorema di Shannon ci dice che se vogliamo trasmettere a una velocità 
inferiore alla capacità del canale, possiamo farlo con una probabilità di errore 
arbitrariamente bassa, utilizzando codici di canale appropriati. 
Se invece vogliamo trasmettere a una velocità superiore alla capacità del canale, 
non importa quanto sofisticati siano i codici di canale che utilizziamo, non saremo 
in grado di ridurre la probabilità di errore al di sotto di un certo livello.

Nel gergo comune diciamo che non abbiamo banda, in realtà non abbiamo potenza del 
segnale, o abbiamo troppo rumore, quindi la capacità del canale è bassa. 

Teorema di Shannon: è possibile trasmettere dati attraverso un canale di comunicazione
affetto da rumore AWGN (Additive White Gaussian Noise) con bitrate R<C.

\paragraph{Rilevamento e correzione degli errori}
le tecniche di rilevazione degli errori sono due principali:
- Automatic repeat request (ARQ): il mittente invia i dati ed anche un codice a rilevazione
d'errore, che sarà utilizzato in ricezione per individuare gli eventuali errori, ed in 
tal caso cheidere la ritrasmissione dei dati corrotti. In modi casi la richiesta è
implicita; il destinatatio invia un acknowledgemebt (ACK) di corretta ricezione dei dati, 
ed il ...
- Forward error correction (FEC): il trasmettitore aggiunge bit di ridondanza al 
messaggio originale, in modo che il ricevitore possa rilevare e correggere gli errori 
senza dover richiedere una ritrasmissione, ad esempio utilizzando codici di correzione 
degli errori come i codici di Hamming o i codici Reed-Solomon.

Come controllo che c'è errore e ridondanza? 
Uso la codifica di hamming, si usa nella codifica di sorgente, quando ho estratto i campioni,
li codifico in bit, e se c'è un errore, la distanza di hamming mi dice quanti bit sono 
diversi tra il messaggio originale e quello ricevuto, se la distanza è 0, non c'è errore, 
se è 1, c'è un errore, se è 2, ci sono due errori, ecc.\\

Il codice di Hamming è un codice di correzione degli errori che può rilevare e 
correggere errori singoli, e rilevare errori doppi.\\

Forward Error Correction (FEC) è una tecnica di codifica che consente al ricevitore 
di correggere gli errori senza dover richiedere una ritrasmissione, utilizzando bit 
di ridondanza aggiunti al messaggio originale.\\

Tutti i sistemi di comunicazione introducono un errore e tutti i sistemi devono usare il BER
(Bit Error Rate) per misurare la qualità del canale, che è il numero di bit errati diviso
per il numero totale di bit trasmessi.\\

Il FEc è il sistema di codifica del canale e aggiunge bit di ridondanza per permettere la 
rilevazione e la correzione degli errori, causati nel canale di comunicazione. 
Ci sono vari tipi di FEC, quella sistematica e quella non sistematica.\\

Codici a blocchi, sono una categoria di algoritmi per la correzione degli errori (FEC) che
operano su segmenti di dati di lunghezza fissa. Il trasmettitore aggiunge bit di controllo 
a un blocco di bit di dati. Il ricevitore usa quei bit per verificare se ci sono stati errori e, 
se possibile, correggerli.

Tipi di codici a blocchi:
- Codici di Hamming: possono correggere errori singoli e rilevare errori doppi.
Il limite di questi codici è che devo mandare tanti bit per correggere pochi errori.
- Codici di Reed-Solomon: possono correggere errori multipli e sono ampiamente 
utilizzati in applicazioni come CD, DVD e comunicazioni satellitari.
Invece di tasmettere la sequenza di bit, considero che i miei simboli sono coefficienti
di un polimonio di grado k-1. Se ho un polimonio di k-1 o lo descrivo con i coeffiecienti 
o con i punti, per esempio x=1 il polimonio passa per un determinato punto.
Lavora prendendo una sequenza di 239 bit, ne aggiunge 16 per trasmettere 255 bit, e 
con 16 bit di controllo, correggo 8 bit.
Si riferisce a polimoni di gradi 8.
- Codici LDPC (Low-Density Parity-Check): sono codici a blocchi che offrono prestazioni 
vicine al limite di Shannon e sono utilizzati in standard di comunicazione come Wi-Fi e 5G.


Controllo degli errori in BLE:

Cyclic Redundancy Check (CRC)
Codici Convoluzionali


\section{Codifica di Linea}
Ci occupiamo della modulazione, e si può chiamare anche codifica di linea, 
tutti i sistemi che utilizziamo dobbiamo permettere agli utenti di accedere al mezzo 
i metodi per permettere la condivisione della risorsa fisica si chiama multiplexing o 
multiple access. 

Una volta che tutti gli utenti hanno trasmesso nel canale di tramessione in ricezione 
facciamo le operazioni inverse, non ci occupiamo più della codifica di sorgente, 
solo richiami della codifica di canale. Adesso entriamo nella parte più importante del 
corso.\\

La codifica di linea deve emeterre i bit supportante fisica, devo mettere il segnale digitale
sulla supportante fisica. Esempio filo di rame la supportante fisica è la tensione.

Se il segnale che voglio trasmettere è in fibra ottica, allora devo trasformare il segnake
in un segnale luminoso, lavoriamo nell'infrarosso.
In ricezione devo fare la demodulazione da ottico a elettrico.

Vediamo come deve essere fatto il segnale per rappresentare la sequenza di 0 e 1, 
on-off keying o non return to zero (NRZ), se voglio trasmettere un 1 accendo la luce, se voglio 
trasmettere uno 0 spengo la luce, è un sistema di modulazione digitale, ma non è molto 
efficiente, perché se ho una lunga sequenza di 0, non trasmetto nulla, e se ho una 
lunga sequenza di 1, trasmetto sempre la luce, quindi non c'è sincronizzazione tra il 
trasmettitore e il ricevitore, non so quando inizia e quando finisce un bit, quindi 
non è molto efficiente.\\

Un altro sistema è il return to zero (RZ), se voglio trasmettere un 1 accendo la luce 
per metà del tempo, e poi spengo la luce, se voglio trasmettere uno 0 spengo la luce per
tutto il tempo, è un sistema di modulazione digitale, ma non è molto efficiente,
perché se ho una lunga sequenza di 0, non trasmetto nulla, e se ho una lunga sequenza di 
1, trasmetto sempre la luce, quindi non c'è sincronizzazione tra il trasmettitore e il
ricevitore, non so quando inizia e quando finisce un bit, quindi non è molto efficiente.\\

Interferenza intersimbolica (ISI): quando il segnale di un simbolo interferisce con il 
segnale di un simbolo adiacente, causando distorsione e rendendo difficile la corretta 
interpretazione dei simboli ricevuti.\\

Criterio di Nyquist: per evitare l'interferenza intersimbolica, la forma d'onda del 
segnale deve essere progettata in modo tale che il segnale di un simbolo sia zero 
nei punti in cui il segnale di un simbolo adiacente è massimo.\\
La condizione di nyquist nel dominio del tempo è che la forma d'onda del segnale deve
 essere zero nei punti in cui il segnale di un simbolo adiacente è massimo, mentre 
 nel dominio della frequenza è che la banda del segnale deve essere limitata a metà
 della frequenza di simbolo.\\

Impulsi di Nyquist: sono forme d'onda che soddisfano il criterio di Nyquist e
 sono utilizzate per trasmettere simboli digitali senza interferenza intersimbolica. 
 Un esempio comune di impulso di Nyquist è l'impulso di sinc, che ha una forma a 
 campana e soddisfa il criterio di Nyquist sia nel dominio del tempo che nel dominio 
 della frequenza.\\

\subsection{Trasmissioni in banda base e banda passante}
La trasmissione in banda base si riferisce alla trasmissione di segnali digitali senza
modulazione, mentre la trasmissione in banda passante si riferisce alla trasmissione di
segnali digitali modulati su una portante a frequenza più alta.\\
La trasmissione in banda base è più semplice e meno costosa, ma è limitata dalla
distanza e dalla qualità del canale, mentre la trasmissione in banda passante 
consente di trasmettere segnali digitali su distanze più lunghe e attraverso canali 
di qualità inferiore, ma è più complessa e costosa.\\
....
....
....
\section{Multiplexing}
Deve permettere a piú utenti di accedere al canale di comunicazione. Questa 
tecnica è chiamata multiplexing o multiple access. Quindi é la tecnica di 
trasmettere insieme più segnali su un unico canale di comunicazione, in modo
che più utenti possano condividere la stessa risorsa fisica.\\

Quasi sempre le operazioni di multiplexing non modificano i segnali, ma li
combinano in modo da poterli trasmettere insieme.\\

Quale tecnologia possiamo usare per multiplexing? 
\begin{itemize}
    \item \textbf{Time Division Multiplexing (TDM)}: divide il tempo in slot e assegna ogni 
    slot a un utente diverso, permettendo a più utenti di trasmettere in momenti diversi.
    E' applicata nella telefonia, dove ogni chiamata viene assegnata a uno slot di tempo specifico, 
    e nella trasmissione dati, dove i pacchetti di dati vengono trasmessi in slot di tempo specifici.
    \item \textbf{Frequency Division Multiplexing (FDM)}: divide la banda del canale in frequenze 
    e assegna ogni frequenza a un utente diverso, permettendo a più utenti di trasmettere 
    contemporaneamente su frequenze diverse. E' applicata nella trasmissione radio, dove ogni 
    stazione radio trasmette su una frequenza specifica, e nella televisione, dove ogni canale 
    televisivo trasmette su una frequenza specifica.
    \item \textbf{Code Division Multiplexing (CDM)}: assegna a ogni utente un codice unico e permette 
    a più utenti di trasmettere contemporaneamente sulla stessa frequenza, distinguendo i
    segnali in base ai codici. E' applicata nelle comunicazioni militari, dove ogni unità militare
    ha un codice unico per comunicare, e nelle reti cellulari, dove ogni utente ha un codice unico per 
    comunicare con la torre cellulare.
\end{itemize}

Non é detto che ogni utente é l'utente finale, potrebbe essere un flusso di dati, ad esempio un video, 
che viene trasmesso da un server a un client, e il server potrebbe essere considerato un utente del 
canale di comunicazione.\\

\subsection{Time Division Multiplexing - TDM}
A ogni slot assento a un utente diverso, e mettendo insieme gli slot ottengo
un frame.\\

Devo trasmettere con una frequenza di campionamento 2 volte la banda del segnale, per trasmettere 
con correttezza, devo trasemttere 8000 byte al secondo.
Tra un segnale e un altro c'è un intervallo di 125 microsecondi, quindi in quel tempo posso trasmettere 
altri utenti, se ho 32 utenti, ogni utente trasmette ogni 4 ms, quindi ogni utente trasmette 250 campioni/s, 
quindi un bitrate di 2 kbps.\\

Se io non devo trasmettere dati quella risorsa viene persa, se un utente pagava di piú gli davo due slot
invece di uno, ma se non trasmetteva nulla, quella risorsa veniva persa, quindi non era efficiente.\\
Anche nelle reti di accesso, l'accesso si basa sul tempo, perché é il modo piú semplice 
di mettere insieme i dati.\\

Hanno inventato il multiplexing statistico, che non assegna slot fissi agli utenti, ma assegna slot 
dinamicamente in base alla domanda di traffico, migliorando l'efficienza del canale di comunicazione.\\

Come Funziona?\\
I segnali vanno memeorizzati e poi affasciati insieme in un frame, e poi trasmessi, in ricezione devo 
fare l'operazione inversa, estrarre i segnali dai frame, e poi decodificarli.\\

Quando la centrale riceve i dati, deve capire a quale utente appartengono, quindi ogni utente ha un identificatore unico,
che viene incluso nei dati trasmessi, in modo che la centrale possa instradare i dati al destinatario corretto.\\

Quando ricevo dalla centrale telefonica, so quale é il mio, so il mio tempo, e butto tutto il resto. 
E'una tecnica molto economica, perché non devo trasmettere dati di controllo, ma è inefficiente, 
perché se un utente non trasmette nulla, quella risorsa viene persa.\\

Confronto con PCM: \\
Nella TDM si é creato la Digital Hierarchy, che é una gerarchia di livelli di multiplexing, 
dove ogni livello combina più flussi di dati in un flusso di dati a velocità più alta.\\
Il livello più basso è il livello 0, che combina 24 canali di voce a 64 kbps ciascuno, per un totale di 1.544 Mbps (T1).\\
Il livello successivo è il livello 1, che combina 4 flussi T1 per un totale di 6.312 Mbps (E1).\\
Il livello successivo è il livello 2, che combina 4 flussi E1 per un totale di 34.368 Mbps (T2).\\
Il livello successivo è il livello 3, che combina 4 flussi T2 per un totale di 139.264 Mbps (T3).\\
Il livello successivo è il livello 4, che combina 4 flussi T3 per un totale di 565.148 Mbps (T4).\\
Il livello più alto è il livello 5, che combina 4 flussi T4 per un totale di 2.260.592 Mbps (T5).\\
 min 37


\section{Orthogonal Frequency Division Multiplexing - OFDM}
Non è un multiplexing, non significa che stiamo dando accesso a più utenti, è una tecnica 
di modulazione del segnale, è così potente che ha rivoluzionato i sistemi di comunicazione, 
è usato nella fibra ottica, è la base dei sistemi wi fi, nelle reti cellulari (4G e 5G), nelle trasmissioni televisive digitali,
ADSL/VDSL in quest'ultimi però si chiama DMT (Discrete Multitone Modulation).\\

OFDM è una tecnica di modulazione che divide un canale di comunicazione in più sottocanali ortogonali,
ciascuno dei quali trasmette una porzione del segnale complessivo, permettendo di trasmettere dati a 
velocità più elevate e migliorando la resistenza alle interferenze e alla distorsione del canale.\\

Quando trasmetto attraverso canale wireless, spesso il canale a causa del multipath fading, 
attenua alcune frequenze più di altre, questo causa distorsione del segnale, ma se allungo il tempo simbolo, 
restringo la banda e quindi lo spettro lo riduco, quindi il canale è più piatto, e quindi è più facile trasmettere, ma se allungo troppo il tempo simbolo,
aumenta la probabilità di interferenza intersimbolica, quindi OFDM divide il canale in più sottocanali (diverso) ortogonali, 
ognuno con una banda più stretta, permettendo di trasmettere dati a velocità più elevate e migliorando la resistenza 
alle interferenze e alla distorsione del canale.\\

Quindi prendo la banda disponibile, per esempio per wi fi ho 20 kHz, e la divido in 64 sottocanali, ognuno con una banda di 
312.5 Hz, e ogni sottocanale trasmette una porzione del segnale complessivo, quindi se ho 64 sottocanali, posso trasmettere 64 
simboli contemporaneamente, e quindi aumentare la velocità di trasmissione.\\

Efficienza spettrale: $\frac {R(bit-rate)}{B(banda)}$\\
L'efficienza spettrale di OFDM è elevata, perché permette di trasmettere dati a velocità più elevate utilizzando la stessa banda, 
migliorando l'efficienza spettrale rispetto ad altre tecniche di modulazione.\\
La sinc si annulla per valori come 1, 2 ecc, quindi l'idea è che se devo trasmettere un altro segnale nello spettro lo modulo
per farlo passare per quei valori, in questo modo non interferisce con il segnale precedente, e quindi posso trasmettere più segnali 
contemporaneamente senza interferenza.\\

La distanza tra i due segnali in frequenza è di $\frac 1 T$. 

Se il mio sistema di trasmissione si comporta strano in una delle frequenze posso rendermene conto e posso facilmente equalizzare.\\

\subsection{Multi-carrier modulation}
Non è un multiplexing, non è neanche una frequency division multiplexing, è una tecnica di modulazione che utilizza più portanti per trasmettere dati,
dobbiamo trasformare un segnale con R bit rate in N sottocanali (flussi) attraverso un Serial to Parallel Converter, e poi ciascun flusso
è trasmesso ad una diversa sottofrequenza (\textbf{subcarrier}).
Due sottocanali sono ortogonali se la loro frequenza è multipla di $\frac 1 T$ (distano tra loro di $\frac 1 T$), dove T è il tempo simbolo, 
in questo modo non interferiscono tra loro.\\

Ci aggiungo a questa sottoportante la portante vera, per esempio se la portante è a 2.4 GHz, e la sottoportante è a 312.5 Hz, allora 
la frequenza del segnale trasmesso sarà 2.4003125 GHz, e così via per gli altri sottocanali.\\

Discrete Multitone Modulation (DMT) è una tecnica di modulazione multi-carrier utilizzata principalmente nelle comunicazioni 
su linee telefoniche, come ADSL e VDSL, per trasmettere dati a velocità elevate su distanze relativamente lunghe. 
DMT divide la banda disponibile in più sottocanali, ognuno dei quali trasmette una porzione del segnale complessivo, 
migliorando l'efficienza spettrale e la resistenza alle interferenze e alla distorsione del canale.\\

Allargo la durata del bit, quindi la durata del simbolo è $\frac N R$, ho ristretto lo spettro. La separazione di ogni sinc 
è fatta in modo da non avere interferenza tra i sottocanali.\\

Come faccio a filtrarle visto che sono sovrapposte? Uso il principio di ortogonalità, i due segnali sono ortogonali se l'integrale del loro 
prodotto nella durata del simbolo è uguale a zero (ovvero delta di Dirac).

Scelgo gli esponenziali complessi come funzioni di base, sono chiamati kernel di Fourier, e come frequenza scelgo $\frac m T$, 
dove m è un intero, in questo modo ottengo la trasformata di Fourier discreta (DFT), che mi permette di trasformare il segnale dal dominio del 
tempo al dominio della frequenza, e viceversa.\\

Quindi se prendo la frequenza $\frac 1 T$ ho un coseno, se prendo la frequenza $\frac 2 T$ ho un coseno con frequenza doppia, e così via, 
in questo modo ottengo i sottocanali ortogonali.\\

Poichè se faccio trasformata di una rect ottengo una sinc, se la rect è moltiplicata per un coseno, ottengo una sinc traslata.
\subsection{Fast Fourier Transform - FFT}
La Fast Fourier Transform (FFT) è un algoritmo efficiente per calcolare la trasformata di Fourier discreta (DFT) di un segnale,
che è una rappresentazione del segnale nel dominio della frequenza.\\

Il processamento del segnale è fatto a livello digitale, ci faccio la trasformata di fourier e poi lo mando al DAC (lo trasforma ad analogico) per trasmetterlo.
S/P e P/S, per trasformare il segnale da seriale a parallelo e viceversa.\\

Le serie / parallelo dividono il flusso di dati (che hanno un bit rate alto) in più flussi di dati (con bit rate più basso) 
che possono essere trasmessi su N sottocanali ortogonali, quindi i vari flussi divisi avranno un bit rate di $\frac R N$, e poi alla fine li ricombino in seriale.\\
Le sottoportante le scelgo con frequenza $\frac 1 T$ perchè è l'inverso del tempo di simbolo.

Tornato da parallelo a seriale aggiungo un prefisso ciclico. Poi nella trasmissione devo fare la modulazione, quindi moltiplico per una portante, e poi trasmetto.\\
In ricezione è uguale ma prima devo demodulare, quindi moltiplico per la portante, poi tolgo il prefisso ciclico, poi faccio S/P, poi la FFT per 
trasformare dal dominio del tempo al dominio della frequenza, e poi faccio il parallelo a seriale per ricombinare i flussi di dati.\\

\paragraph{Prefisso ciclico}
Sto trasmettendo il segnale come flusso di bit, a causa della distorsione, questi bit si possono sovrapporre, causando l'interferenza intersimbolica, 
OFDM ha pensato che invece d mettere uno spazio vuoto tra un simbolo e l'altro, metto una copia del simbolo precedente (una parte del simbolo), 
in questo modo se c'è interferenza intersimbolica, il simbolo precedente è uguale a quello che sta arrivando, quindi non c'è interferenza, 
e posso ricostruire il segnale originale.\\

Più lungo è il prefisso ciclico, più è probabile che riesca a eliminare l'interferenza intersimbolica, ma più è lungo il tempo di simbolo, quindi
più è bassa l'efficienza spettrale, quindi c'è un tradeoff tra la lunghezza del prefisso ciclico e l'efficienza spettrale. \\%daje

Questo mi serve nei canali wireless per disperdere il multipath fading, nella fibra ottica per la dispersione cromatica, 
e in generale per qualsiasi canale che introduce distorsione.\\

\subsection{Schema del trasmettitore OFDM}
Il segnale di partenza lo divido in più flussi con convertitore S/P, in modo che ogni flusso abbia un Bit rate che è il bit rate di partenza diviso 
per il numero di flussi, creo una costellazione, ovvero 4 simboli e ogni simbolo rappresenta 2 bit (se ho 16 simboli ogni simbolo rappresenta 4 bit), 
OFDM poi fa il FFT per trasformare dal dominio del tempo al dominio della frequenza, poi P/S, poi aggiungo il prefisso ciclico, poi DAC,  
poi moltiplico per la portante (parte reale usa coseno, parte immaginaria usa coseno), e poi trasmetto.\\

In ricezione, prima moltiplico per la portante, lo filtro quindi tolgo la portante (rimetto in banda base), poi ADC, poi tolgo il prefisso ciclico, 
poi S/P, poi faccio la FFT per trasformare dal dominio del tempo al dominio della frequenza, poi faccio il parallelo a seriale per ricombinare i flussi di dati, 
e poi decodifico i simboli per ottenere i bit.\\

\subsection{Prestaziomi dell'OFDM}
PRO:
\begin{itemize}
    \item Robustezza rispetto all'interferenza intercanale, all'interferenza intersimbolica e al multipath fading, 
    grazie alla divisione del canale in sottocanali ortogonali e all'uso del prefisso ciclico.
    \item Efficienza spettrale elevata, grazie alla possibilità di trasmettere dati a velocità più elevate utilizzando la stessa banda.
    \item Permette formati di modulazione avanzati
    \item Implementazione efficiente, grazie all'uso della FFT per la modulazione e demodulazione dei sottocanali.
    \item Bassa sensibilità agli errori di sincronizzazione, grazie alla divisione del canale in sottocanali ortogonali e all'uso del prefisso ciclico.
    \item Processamento elettronico del segnale
\end{itemize}
CONTRO: 
\begin{itemize}
    \item Sensibile all'effetto doppler
    \item Sensibile ai problemi di sincronizzazione in frequenza
    \item Altro peak-to-average power ratio (PAPR), che può causare problemi di amplificazione del segnale e ridurre l'efficienza energetica del sistema.
    \item Perdita di efficienza dovuta all'inserimento del ciclo prefisso, che riduce la quantità di dati trasmessi per unità di tempo.
\end{itemize}

\section{Complementary Cumulative Distribution Function - CCDF}
La Complementary Cumulative Distribution Function (CCDF) è una funzione che descrive la probabilità che una variabile casuale superi un certo valore, 
ed è spesso utilizzata per analizzare il Peak-to-Average Power Ratio (PAPR) nei sistemi di comunicazione, come OFDM, per valutare la probabilità che 
il PAPR superi un certo livello, aiutando a progettare sistemi più efficienti e robusti contro le distorsioni causate da picchi di potenza elevati.\\

Peak to Average Power Ratio (PAPR): misura per quanto tempo il segnale che devo trasmettere supera il valore medio. 
CCDF è la probabilità che il PAPR superi un certo livello, ad esempio se voglio sapere la probabilità che il PAPR superi 10 dB, guardo la CCDF a 10 dB.\\

Il fatto che la dinamica del mio segnale è così ambia è un problema, si possono creare effetti non lineare, perché la potenza è troppo elevata e la risposta 
della fibra ottica non è lineare.\\

Questa tecnica la usiamo nel wi fi, in ottica e nei sistemi mobili.

\section{Discrete Multi Tone Modulation - DMT}
E' l'OFDM in banda banse e viene utilizzato nell'ADSL.
Nel 1960 abbiamo avuto i primi sistemi in rame, funzionava con un sistema analagico e quelle
rete di chiamavano POTS (Plain Old Telephone Service), si chiamavano old proprio perchè ora 
non esistono più, erano reti di comunicazione analogiche, e quindi non erano digitali, e quindi non potevano trasmettere dati, ma solo voce.\\

Il segnale voce era trasformato in segnale elettrico e trasmesso attraverso il filo di rame, e poi in ricezione veniva trasformato di nuovo in segnale acustico.\\

I primi 4 kHz non li tocchiamo (sono il segnale voce), dividiamo in sottoportanti la banda restante.

ADSL (Asymmetric Digital Subscriber Line) è una tecnologia di comunicazione che utilizza le linee telefoniche in rame 
per trasmettere dati a velocità elevate, e DMT è la tecnica di modulazione utilizzata in ADSL per dividere la banda 
disponibile in più sottocanali, ognuno dei quali trasmette una porzione del segnale complessivo, migliorando l'efficienza spettrale e l
a resistenza alle interferenze e alla distorsione del canale.\\

Considero 512 sottoportanti, ognuna con una banda di 4.3125 kHz, e ogni sottoportante trasmette una porzione del segnale complessivo, 
quindi se ho 512 sottoportanti, posso trasmettere 512 simboli contemporaneamente, e quindi aumentare la velocità di trasmissione.\\

Applico a tutti questi i simboli la trasformata di fourier, e poi li trasmetto, in ricezione faccio l'operazione inversa, 
quindi faccio la trasformata di fourier inversa, per trasformare dal dominio della frequenza al dominio del tempo, 
e poi decodifico i simboli per ottenere i bit.\\

Bit loading: il sistema misura i dati in trasmissione se c'è rumore o interferenza (che chiamiamo diafonia), 
e diminuisco il numero di bit che trasmetto (uso Shannon). Più c'è rumore, meno bit trasmetto, e viceversa.\\

Poi applico il QAM (Quadrature Amplitude Modulation) per modulare i simboli, poi trasformata di Fourier, poi +CP e poi DAC.\\

VDSL2 è quella che uso se a casa mia non arriva la fibra. Il tratto di rame deve essere più corto possibile, 
più è lungo più c'è il rumore, quindi più è basso il bit rate, e più è corto più è alto il bit rate.\\

I sistemi in banda base sono quelli che viaggiano in rame. 

\chapter{Sistemi Wireline e Wireless}
\section{Sistemi Wireless vs Wireline}

Vantaggi Wireline:
\begin{itemize}
    \item Maggiore velocità di trasmissione
    \item Maggiore affidabilità e stabilità
    \item Sicurezza più elevata contro le intercettazioni
\end{itemize}
Svantggi Wireline:
\begin{itemize}
    \item Limitata mobilità per i dispositivi
    \item Infrastruttura fissa e costosa da installare
\end{itemize}
Vantaggi Wireless:
\begin{itemize}
    \item Maggiore mobilità e flessibilità
    \item Scalabilità più semplice per nuovi dispositivi
\end{itemize}
Svantaggi Wireless:
\begin{itemize}
    \item Suscetibilià a interferenze e perdite di segnale
    \item Maggiore rischio di sicurezza e intercettazioni
    \item Latenza e variabilità della velocità di trasmissione
\end{itemize}

\section{Cavi in rame e ottici}
I cavi in rame studiamo il doppino telefonico, che è una coppia di fili intrecciato, 
Cavo ethernet, che sono 4 coppie di fili intrecciati, e il cavo coassiale, che è un filo centrale circondato da un isolante e una schermatura.
I cavi in fibra ottica invece utilizzano la luce per trasmettere i dati, e sono composti da un nucleo di vetro o plastica circondato da un rivestimento riflettente e una guaina protettiva.\\

Nella fibra ottica (che è un vetro) non ci viaggiano elettroni ma impulsi luminosi, ci sarà una componente
che converte da elettrico a luminoso.

\subsection{Doppino di Rame}
E' composto da due conduttori di rame (di diametri leggermente diversi) intrecciati tra loro, questo intreccio si chiama binatura, 
e ogni filo di rame ha una guaina isolante. La binatura serve a ridurre le interferenze elettromagnetiche, e permette di trasmettere dati a velocità fino a 100 Mbps su distanze fino a 100 metri.

A riposo sul filo c'è ugualmente una differenza di potenziale costante. Sul doppino telefonico viaggia
il POTS, l'ISDN, lADSL, VDSL, e G.Fast.\\

Per prima cosa si è passati dal POTS (Plain Old Telephone Service) all'ISDN (Integrated Servuce Digital Network), poi ADSL con 256 sottoportanti e modulazione QAM.
Con VDSL si è arrivati a 4096 sottoportanti e modulazione QAM-256, e infine con G.Fast si è arrivati a 10624 sottoportanti e modulazione QAM-4096, con velocità fino a 1 Gbps su distanze fino a 100 metri.\\

La tecnica del G Fast, che ha permesso di estendere la banda del rame, ha permesso di estendere fino a 212 MHz.

Il G Fast utilizza una tecnica chiamata vectoring, che consiste nell'utilizzare un algoritmo di cancellazione del rumore per ridurre 
le interferenze tra i cavi, e permette di aumentare la velocità di trasmissione su distanze più lunghe.\\

Utilizza il DFT-Spread OFDM, che è una tecnica di modulazione che combina la modulazione OFDM con la trasformata discreta di Fourier (DFT) (la applica due volte), e permette di ridurre 
l'effetto del fading e migliorare la resistenza alle interferenze.\\

Quando si hanno più doppini ritorti è importante che la twistatura dei fili sia diversa tra i cavi, per evitare interferenze tra di essi (evitano l'effetto di diafonia).\\

Tecnica di trasmissione differenziale: si trasmette il segnale su entrambi i fili del doppino, e si utilizza la differenza di potenziale tra i due fili per trasmettere i dati, 
questo permette di ridurre le interferenze elettromagnetiche e migliorare la qualità del segnale.\\

Il connettore del doppino telefonico è il connettore RJ11, che ha 4 o 6 pin, e viene utilizzato per collegare i telefoni alla rete telefonica.\\
Il RJ10 è quello che connette i telefoni alla rete interna, e ha 4 pin.\\

\subsection{Cavo Ethernet}
Il cavo Ethernet è composto da 4 coppie di fili intrecciati, e viene utilizzato per collegare i dispositivi alla rete locale (LAN). 
Esistono diverse categorie di cavi Ethernet, che si differenziano per la velocità di trasmissione e la qualità del segnale.\\

Ci sono i cavi standard che si chiamano Unshielded Twisted Pair (UTP), che sono i più comuni, e i cavi Shielded Twisted Pair (STP), che hanno una schermatura aggiuntiva per ridurre le interferenze elettromagnetiche.\\
I Foiled Twisted Pair (FTP) hanno una schermatura globale, e i Screened Twisted Pair (ScTP) ogni coppia (doppino) viene schermato.\\

RJ-45 è il connettore standard per i cavi Ethernet, e ha 8 pin, e viene utilizzato per collegare i dispositivi alla rete locale (LAN).\\

Due modi per connetterli: cross e patch.
Quando faccio una connessione diretta è patch, mentre uso cross se devo connettere pc con pc o switch ethernet con altro switch ethernet.\\

Vantaggi:
\begin{itemize}
    \item Alta Affidabilità e stabilità
    \item Standard Universale
    \item Supporto PoE (Power over Ethernet) per alimentare dispositivi
    \item Costo contenuto per distanze fino a 100 metri
    \item Ampia compatibilità con dispositivi di rete
\end{itemize}

Svantaggi:
\begin{itemize}
    \item Limitato a 100 metri senza ripetitori
    \item Suscettibile a interferenze elettromagnetiche in ambienti rumorosi
    \item velocità elevate richiedono Cat 6A o superiore
\end{itemize}

\subsection{Cavo Coassiale}
Il cavo coassiale è composto da un filo conduttore di rame (anima) e un dielettrico (polietilene) garantisce 
l'isolamento tra il centro del conduttore e uno schermo di metallo intrecciato (maglia).

Le distanze sono limitate perché il segnale si attenua, e la velocità di trasmissione dipende dalla qualità del cavo e dalla distanza, ma può raggiungere fino a 10 Gbps su distanze fino a 100 metri.\\
Noi li usiamo per la televisione via cavo, e per la connessione a Internet tramite modem via cavo.\\

\section{Spettro elettromagnetico}
Lambda per convenzione è la lunghezza d'onda (distanza tra un picco e un altro), quando un campo magnetico si propaga coper
uno spazio e la distanza tra un picco e l'altro è la lunghezza d'onda, e la frequenza è il numero di cicli al secondo (Hz).\\

Si propaga con una velocità c che è la velocità della luce, e $c = \lambda * f.$\\
Distinguiamo i campi elettromagnetici in base alla frequenza, e abbiamo le onde radio (fino a 300 GHz, frequenze basse), le microonde (da 300 MHz a 300 GHz), 
l'infrarosso (da 300 GHz a 400 THz) le lunghezze onda che utilizziamo nella fibra ottica, la luce visibile (da 400 THz a 800 THz), 
l'ultravioletto (da 800 THz a 30 PHz), i raggi X (da 30 PHz a 30 EHz) e i raggi gamma (oltre 30 EHz).\\

Utilizziamo le micro onde per le comunicazioni wireless, e le onde radio per le comunicazioni a lunga distanza.\\

\subsection{Proprietà della luce}
Riflessione: quando un'onda elettromagnetica colpisce una superficie, parte dell'energia viene riflessa, e la quantità di energia riflessa dipende dalle proprietà della superficie e dall'angolo di incidenza.\\
Rifrazione: quando un'onda elettromagnetica passa da un mezzo a un altro, la sua velocità e direzione cambiano a causa della differenza di indice di rifrazione tra i due mezzi.\\
Diffrazione: quando un'onda elettromagnetica incontra un ostacolo o una fessura, si piega e si diffonde intorno ad esso, e la quantità di diffrazione dipende dalla lunghezza d'onda e dalle dimensioni dell'ostacolo o della fessura.\\
Diffusione (Scattering): quando un'onda elettromagnetica interagisce con particelle o irregolarità nell'ambiente, viene dispersa in diverse direzioni, e la quantità di diffusione dipende dalle proprietà delle particelle e dalla lunghezza d'onda dell'onda elettromagnetica.\\
Polarizzazione: la direzione del campo elettrico di un'onda elettromagnetica può essere orientata in diverse direzioni, e la polarizzazione può essere lineare, circolare o ellittica.\\
Interferenza: quando due o più onde elettromagnetiche si sovrappongono, possono interferire costruttivamente (aumentando l'ampiezza) o distruttivamente (riducendo l'ampiezza), e la quantità di interferenza dipende dalle fasi relative delle onde.\\
\section{Fibra ottica}

Dati sono trasmessi come impulsi di luce, e la fibra ottica è composta da un nucleo di vetro o plastica circondato da un rivestimento riflettente e una guaina protettiva.\\
La rete telefonica Plain Old Telephone Service (POTS) è stata la prima tecnologia di comunicazione a utilizzare la fibra ottica, e ha permesso di trasmettere dati a velocità fino a 1
Gbps su distanze fino a 100 km.\\

Negli anni 90 sono nati i modem, prendevano il segnale digitale e lo rendevano analogico per poterlo trasmettere sulla rete telefonica, e poi lo riconvertivano in digitale all'arrivo.\\
Successivamente ISDN ha permesso di trasmettere dati digitali direttamente sulla rete telefonica, e ha permesso di raggiungere velocità fino a 128 Kbps.\\

I sistemi su rame è punto punto, non hanno inventato niente per far accedere più persone sullo stesso doppino, perché all'epoca c'era tanto rame e meno persone da dover connettere. 
Rete parte dalla casa del cliente e va all'armadio di distribuzione, e poi arriva alla centrale telefonica, e da lì va alla rete di trasporto.\\

C'è un sistema di cavi che poi si separano per raggiungere le case dei clienti.\\
Come operatore un sistema punto punto è difficile e quindi hanno inventato il sistema fibra ottica, che è un sistema punto multipunto, e permette di condividere la fibra ottica tra più clienti.
Le centrali venivano chiamate centrali di linea, e ora si chiamano armadi di distribuzione, e sono più vicine ai clienti, e permettono di condividere la fibra ottica tra più clienti.\\

\subsection{Protocolli di trasmissione}
\paragraph{ADSL - Asymmetric Digital Subscriber Line}
E' una tecnologia di trasmissione dati che utilizza la rete telefonica in rame per fornire connessioni a banda larga.
Il primo ADSL ha il doppino telefonico con una banda 1,1 MHz, però per trasmettere non possiamo farlo con cifre elevate, quindi suddividiamo la banda in sottocanali. \\

La voce era il primo bin, poi dovevo filtrarla quindi i 4 bin erano vuoti, poi (quelli verdi) upstream, e quelli rossi downstream.\\

Nel rame più aumento il cammino più aumento il cross talk, e quindi più aumenta il rumore, e più diminuisce la velocità di trasmissione.\\

Su ogni sottocanale posso trasmettere dei simboli (non solo 0 e 1) se ho un SNR sufficientemente alto, posso trasmettere su ogni sottocanale un numero di bit che dipende dal SNR, e quindi 
posso aumentare la velocità di trasmissione.\\

Posso caricare su ogni sottoportante un massimo di 15 bit, ma in media non si superano 8-10 bit. Siccome ogni sottoportante ha una larghezza di 4.3 kHz, l'inverso della larghezza della sottoportante 
è il tempo di simbolo, ovvero 0.23 ms, poi aggiungo un prefisso ciclico di 0.03 ms, e quindi il tempo totale di simbolo è di 0.26 ms, 
e quindi la velocità di trasmissione è di 1/0.26 ms = 3.85 k simboli al secondo (approssimiamo a 4Kbaund). Se ogni sottoportante trasmette più simboli (più bit contemporaneamente)
posso aumentare la velocità di trasmissione, e quindi posso raggiungere velocità fino a 24 Mbps in downstream e 1.4 Mbps in upstream.\\

Con ADSL2+ si è arrivati a 2.2 MHz di banda, e quindi a velocità fino a 24 Mbps in downstream e 1.4 Mbps in upstream su distanze fino a 1.5 km.\\

DSLAM (Digital Subscriber Line Access Multiplexer) è un dispositivo che si trova nella centrale telefonica, e permette di aggregare le connessioni ADSL dei clienti e di 
collegarle alla rete di trasporto.\\

Tutta la rete fino ai cabinet di distribuzione ora è in fibra ottica, e poi dal cabinet di distribuzione alla casa del cliente è in rame, e questo sistema si chiama Fiber to the Cabinet (FTTC).
In genere si usa o 17a (banda rame è 17MHz) o 35b (banda rame è 35MHz), quindi riusciamo a dare banda fino a 300 Mbps in downstream e 100 Mbps in upstream su distanze fino a 500 metri.\\

I sistemi attuali sono Fiber to the Home (FTTH), dove la fibra ottica arriva direttamente a casa del cliente, e permette di raggiungere velocità fino a 1 Gbps in downstream e 500 Mbps in upstream su distanze fino a 20 km.\\
Il G.fast estende le prestazione della fibra usando l'ultimo tratto in rame.

\paragraph{Spettro elettromagnetico}
Lavoriamo con la banda dello spettro elettromagnetico che non è visibile, lavoriamo con il Near Infrared. Mentre le radiofrequenze sono utilizzate per le comunicazioni wireless, e le microonde 
per le comunicazioni a lunga distanza. Con il 5G usiamo le frequenze che prima erano utilizzate per la televisione (onde radio).
Quindi per wi fi e bluetooth usiamo le frequenze delle microonde, e per la fibra ottica usiamo le frequenze del Near Infrared.\\

\paragraph{Come è costituita la Fibra Ottica}
E' composta da un core (nucleo) di diamentro 9 micron (standard) o 50 micron, e da un cladding (rivestimento) di diametro 125 micron, e da una guaina protettiva.\\
Rifressione totale, è il fenomeno che permette alla luce di rimanere confinata all'interno del core, e si verifica quando la luce colpisce l'interfaccia tra il core e il cladding con 
un angolo maggiore dell'angolo critico.\\ 
La luce si propaga all'interno del core grazie alla riflessione totale. Quindi all'interno del core la luce si propaga, se la fibra è un po' piegata, subisce la riflessione totale (agisce quasi come specchio),
ho due strati di vetro con indice di rifrazione diverso (core ng, cladding n1). E' necessario che la luce stia in un mezzo con indice di rifrazione più alto rispetto al mezzo circostante, 
quindi ng>n1, e questo permette alla luce di rimanere confinata all'interno del core.\\

L'angolo critico è l'angolo minimo di incidenza per cui si verifica la riflessione totale, e dipende dagli indici di rifrazione del core e del cladding, e può essere calcolato utilizzando 
la legge di Snell.\\

Se la potenza luminosa si perde nel mantello (cladding) o nella guaina protettiva, si verifica un'attenuazione del segnale, e la quantità di attenuazione dipende dalle proprietà del materiale 
e dalla lunghezza d'onda della luce.\\

Apertura numerica (NA) è una misura della capacità di una fibra ottica di raccogliere e trasmettere la luce, e dipende dagli indici di rifrazione del core e del cladding, e può essere calcolata 
utilizzando la formula $NA = \sqrt{ng^2 - n1^2}$.\\

\paragraph{Modalità di propagazione}
Si dicono fibre multimodali quelle fibre ottiche che hanno un core di diametro maggiore di 50 micron, e permettono alla luce di propagarsi in più modalità (percorsi) all'interno del core, e sono
utilizzate principalmente per le comunicazioni a breve distanza (fino a 2 km)\\
Si dicono fibre monomodali quelle fibre ottiche che hanno un core di diametro di 9 micron, e permettono alla luce di propagarsi in una sola modalità (percorso) all'interno del core, 
e sono utilizzate principalmente per le comunicazioni a lunga distanza (oltre 2 km).\\

Perché usare fibra multimodale?
La fibra multimodo soffre dell'effetto dispersione intermodale, ma viene usata nelle data center. Se all'interno della fibra si propagano più modalità, ci sarà un modo che si propaga in modo lineare, 
altri seguendo determinati angoli, quindi la luce in ingresso se eccita più modi si allarga, se la fibra è lunga la luce si allarga così tanto che non riesco più a distinguere i simboli, 
e quindi la velocità di trasmissione diminuisce. Questo effetto si chiama allargamento dell'impulso e può portare alla ISI (Inter Symbol Interference). L'allargamento dell'impulso dipende dalla 
lunghezza della fibra (più lunga fibra più basso bit rate per dispersione intermodale). Un vantaggio della fibra multimodale è che è più economica da produrre e da installare rispetto alla fibra 
monomodale, e può essere utilizzata per le comunicazioni a breve distanza, come ad esempio all'interno di un edificio o di un campus universitario.\\

Ci sono degli standard per le fibre multimodali, come ad esempio OM1, OM2, OM3 e OM4, che si differenziano per la loro capacità di trasmettere dati a diverse velocità e distanze.\\

Quella monomodale invece è utilizzata per le comunicazioni a lunga distanza, come ad esempio tra città o tra continenti, e permette di raggiungere velocità di trasmissione molto elevate 
(fino a 100 Gbps) su distanze fino a 100 km, ma è più costosa da produrre e da installare rispetto alla fibra multimodale.\\
Ha l'effetto chiamato dispersione cromatica, che è causato dalla differenza di velocità di propagazione della luce a diverse lunghezze d'onda, e può portare alla ISI (Inter Symbol Interference).\\
Più aumento la lunghezza della fibra, più gli impulsi arrivano con tempi diversi, e quindi si sovrappongono, e quindi non riesco più a distinguere i simboli, e quindi la velocità di trasmissione diminuisce.\\

Anche nelle fibre multimodali c'è la dispersione cromatica, ma è meno impattante rispetto alla dispersione intermodale, quindi non si nota troppo.\\

Poichè la dispersione cromatica è un effetto lineare noi riusciamo a compensarla, e quindi possiamo raggiungere velocità di trasmissione più elevate su distanze più lunghe. Per farlo 
utilizziamo dei dispositivi chiamati compensatori di dispersione, che possono essere basati su fibre ottiche speciali o su dispositivi a cristalli liquidi, e permettono di compensare la dispersione cromatica
e migliorare la qualità del segnale.\\

ISI per la dispersione intermodale e cromatica.\\

\paragraph{Polarization Division Multiplexing - PMD}
All'interno della singola fibra mando due polarizzazioni diverse (ortogonali), che riesco a separare con un Polarization Beam Splitter. Si usa questa caratteristica nei collegamenti sottomarini.

\paragraph{Attenuazione in una fibra ottica}
La fibra attenua di 0.2 dB/km a 1550 nm, quindi poco. Quando i fotoni si propagano all'interno della fibra, possono imbattersi in impurità o difetti nel materiale, causando così 
la loro dispersione, oppure se la fibra è piegata troppo, la luce può fuoriuscire dal core, causando così un'ulteriore attenuazione del segnale.\\

Le fibre standard attenuano di meno in alcune bande di lunghezza d'onda, quindi si è lavorato in quelle bande. 
Le finestre sono 850 nm, 1300 nm e 1550 nm, e sono state scelte perché in queste bande l'attenuazione è minima, e permettono di raggiungere velocità di trasmissione più elevate su distanze più 
lunghe.\\
La prima finestra (850 nm) non è tanto bassa in attenuazione (2.5 dB/km), e quindi è utilizzata principalmente per le fibre multimodali, mentre la seconda (1300 nm ovvero 1.3 micron) e la terza (1550 nm ovvero 1.55 micron) sono utilizzate 
principalmente per le fibre monomodali.\\

L'attenuazione in fibra è molto bassa rispetto ai cavi in rame, e permette di trasmettere dati a velocità elevate su distanze più lunghe senza la necessità di amplificatori o ripetitori.\\

Quando mando segnale da casa a centrale telefonica uso una banda di 1300 nm (seconda finestra), mentre da centrale a casa 1550 nm (terza finestra), perché in questo modo posso sfruttare al meglio le caratteristiche 
della fibra ottica e raggiungere velocità di trasmissione più elevate su distanze più lunghe.\\
\paragraph{Rumore in una fibra ottica}
Il rumore non è un porblema della fibra ottica, ma è un problema dei dispositivi di trasmissione e ricezione.
Viene introdotto in amplificazione o ricezione, quindi tutte le volte il segnale non è più nel dominio ottico (in realtà ci sono amplificatori ottici) viene introdotto rumore.\\

\paragraph{Connettori}
I connettori per le fibre ottiche sono utilizzati per collegare i cavi in fibra ottica ai dispositivi di trasmissione e ricezione, e sono progettati per garantire una connessione stabile e affidabile.\\
A casa quando portano la fibra fanno un giunto a fusione, devono necessariamente fare una connessione permanente, così continuo (saldo) la fibra. La fibra è di vetro quindi basta riscaldare
un po' la fibra e farla fondere, e poi allineare i due pezzi di fibra in modo che la luce possa passare attraverso di essi senza perdite significative.\\

\paragraph{Wavelenght Division Multiplexing - WDM}
Creo tanti flussi di segnale a diverse lunghezze donda, e li sommo insieme, e poi alla ricezione li separo con un filtro ottico.\\
E' una tecnica di multiplexing che consente di trasmettere più segnali ottici per aumentare la banda disponibile su una singola fibra ottica, e viene utilizzata principalmente nelle comunicazioni a lunga distanza, 
come ad esempio tra città o tra continenti.\\

Nei data center uso 4 lunghezze d'onda in prima finestra chiamata short wavelength division multiplexing (SWDM), e posso raggiungere velocità di trasmissione fino a 100 Gbps su distanze fino a 100 metri.
Nei colleghi Fiber to the Home (FTTH) uso 4 lunghezze d'onda in seconda finestra chiamata coarse wavelength division multiplexing (CWDM), e posso raggiungere velocità di trasmissione fino a 10 Gbps su distanze fino a 20 km.\\
ONU (Optical Network Unit) è un dispositivo che si trova presso il cliente, e permette di convertire il segnale ottico in segnale elettrico, e viceversa, e di collegare i dispositivi del cliente alla rete di trasporto.\\

Nelle reti dorsali dove usiamo il protocollo OTN (Optical Transport Network) separo le lunghezze d'onda usando demultiplexer e le amplifico usando amplificatori ottici, e poi le riassemblo usando multiplexer.
Mi conviene mantenere il segnale ottico il più possibile e non convertirlo in elettrico, perché ogni volta che lo converto introduco rumore, e quindi mantengo il segnale ottico il più possibile per mantenere la qualità del segnale.\\

\paragraph{Dense Wavelength Division Multiplexing - DWDM}
Hanno inventato gli amplificatori ottici che mantengono il segnale ottico, amplificano il segnale e tutte le lunghezze d'onda che sto trasmettendo. 
Quindi ho bisogno per ogni lunghezza d'onda un laser, poi un multiplexer per sommare tutte le lunghezze d'onda, e poi un amplificatore ottico per amplificare il segnale, 
e poi un demultiplexer per separare le lunghezze d'onda alla ricezione.\\

\subsection{Infrastrutture in fibra ottica}
Usiamo la fibra ottica ovunque, all'interno dei data center (Storage Area Network - SAN), nelle reti Backhaul, nelle reti Fiber to the Home (FTTH) e nelle reti sottomarine Backbone(Submarine Cable Network).\\

\paragraph{Data Center}

All'interno dei data center usiamo delle interconnessioni in fibra ottica multidomale per collegare i server tra di loro, e permettere la trasmissione di grandi quantità di dati a velocità elevate.
Bisogna creare connessioni anche da 50 G per ogni singolo server, e quindi si utilizzano cavi in fibra ottica multimodale con connettori LC (Lucent Connector) o MPO (Multi-fiber Push On).\\
Usiamo LED o VCSEL (Vertical Cavity Surface Emitting Laser) come sorgenti di luce, e fotodiodi come rivelatori di luce, e utilizziamo protocolli di trasmissione come ad esempio Ethernet o InfiniBand per gestire la comunicazione tra i server.\\

\paragraph{Reti Backhaul}
Ora con il 5G abbiamo le AAU (Active Antenna Unit) che sono le antenne che si trovano sui tetti, e che permettono di trasmettere i dati tra i dispositivi mobili e la rete di trasporto, 
e sono collegate alla rete di trasporto tramite cavi in fibra ottica.\\

C'è un protoccolo che converte il segnale digitale per trasmetterlo in fibra ottica, e si chiama CPRI (Common Public Radio Interface), e permette di trasmettere dati a velocità fino a 10 Gbps su distanze fino a 20 km.\\

I dati che vengono dalla base transceiver station BTS poi vengono divisi in più flussi di dati, e poi vengono trasmessi in fibra ottica fino alla centrale telefonica, e poi vengono aggregati e trasmessi alla rete di trasporto.\\

\paragraph{Dorsali e Sottomarine}

Le reti dorsali sono le reti che collegano le città tra di loro, e sono costituite da cavi in fibra ottica che permettono di trasmettere dati a velocità elevate su distanze più lunghe.\\
Le reti sottomarine sono le reti che collegano i continenti tra di loro, e sono costituite da cavi in fibra ottica che permettono di trasmettere dati a velocità elevate su distanze più lunghe, e sono progettate per resistere alle condizioni ambientali estreme del fondo oceanico.\\

Per quelli sottomarini gli amplificatori sono ottici ma hanno bisogno di energia elettrica (di differenza di potenziale) quindi oltre al segnalpe viaggia anche l'energia elettrica, e quindi hanno bisogno di stazioni di alimentazione lungo il percorso per fornire energia agli amplificatori ottici.
Il mare alla fibra (che è vetro) non dovrebbe dare problemi ma visto che ci sono altri disturbi come ad esempio le correnti marine, le vibrazioni, le variazioni di temperatura e gli animali è necessario utilizzare cavi in fibra ottica progettati per resistere a queste condizioni ambientali estreme.\\

Nelle reti metropolitane le fibre si inseriscono scavando prima il condotto e poi vengono soffiate al loro interno.\\

\subsection{Evoluzione delle reti in fibra ottica}

\subsection{Collegamento in Fibra ottica}

Quando si lavora ottico si lavora con un fotone, ma le informazioni provengono da elettrico quindi c'è la sorgente che converte da elettrico a ottico. 
Il laser o led emette un fotone che viaggia all'interno della fibra ottica, e poi arriva al rivelatore che converte il segnale ottico in segnale elettrico, e permette di recuperare le informazioni trasmesse.\\

Ho la mia sequenza di dati (bit) io la trasformo in una serie di impulsi elettrici che modula il laser, ovvero accende o spegne il laser, e quindi creo una sequenza di impulsi luminosi che viaggiano all'interno della fibra ottica, 
e poi arrivano al rivelatore che converte il segnale ottico in segnale elettrico, e permette di recuperare le informazioni trasmesse.\\

Filtro ottico è un dispositivo che permette di selezionare una specifica lunghezza d'onda della luce, e di bloccare le altre, così si elimina il rumore e si migliora la qualità del segnale.\\

Come sorgente possiamo usare o i laser o i LED (non vengono usati spesso perché sono più costosi e meno efficienti, quindi non vanno bene per le comunicazioni in fibra ottica).
Se accendo il laser trasmetto il bit 1, se lo spengo trasmetto il bit 0, e quindi creo una sequenza di impulsi luminosi che viaggiano all'interno della fibra ottica, e poi arrivano al rivelatore che converte il segnale ottico in segnale elettrico, e permette di recuperare le informazioni trasmesse.
Questo tipo di modulazione si chiama On-Off Keying (OOK), ed è la forma più semplice di modulazione digitale, e viene utilizzata principalmente nelle comunicazioni in fibra ottica a bassa velocità.\\

Se io voglio accendere un laser a una velocitàa di 10 miliardi al secondo non ci riesco, quindi ho bisogno di un dispositivo chiamato modulatore, che fa passare o meno la luce, sono chiamati
modulatori esterni.\\

\subsection{Sorgente della fibra ottica}
La sorgente è un laser. Per parlare del laser bisogna spiegare cosa è una luce.\\
E' una onda elettromagnetica che si propaga nel vuoto con la velocià della luce, ed è solo una parte dello spettro magnetico che va dai 400 nm (violetto) ai 700 nm (rosso).\\

Plank aveva ipotizzato che la luce fosse composta da particelle chiamate fotoni, e che ogni fotone avesse un'energia quantizzata, ovvero che potesse assumere solo valori discreti di energia.
Più sale la frequenza più sale l'energia del fotone, e quindi più è difficile da generare.\\

Einstein ha dimostrato che la luce è composta da fotoni, e che ogni fotone ha un'energia quantizzata, e che la luce può essere emessa o assorbita solo in pacchetti discreti di energia, chiamati quanti.\\

Il fotone è un onda elettromagnetica che si propaga nel vuoto con la velocità della luce, e ha un'energia quantizzata, che dipende dalla frequenza della luce, 
e può essere calcolata utilizzando la formula $E = h\nu$, dove h è la costante di Planck e $\nu$ è la frequenza della luce.\\

Come interagisce con la materia la luce, 3 modi:
\begin{itemize}
    \item Assorbimento: quando un fotone colpisce un atomo o una molecola, può essere assorbito, e l'energia del fotone viene trasferita all'atomo o alla molecola, 
    e questo processo è chiamato assorbimento, e può essere utilizzato per generare luce coerente, come ad esempio nei laser.\\
    \item Emissione spontanea: quando un atomo o una molecola assorbe energia, può emettere un fotone in modo casuale, e questo processo è chiamato emissione spontanea, 
    e può essere utilizzato per generare luce incoerente, come ad esempio nei LED. LED sta per "Light Emitting Diode".\\
    \item Emissione stimolata: quando un fotone colpisce un atomo o una molecola che è già eccitato, se è colpito da un altro fotone, può stimolare l'emissione di un altro fotone con la stessa frequenza, fase e direzione, 
    e questo processo è chiamato emissione stimolata, e può essere utilizzato per generare luce coerente, come ad esempio nei laser. LASER sta per "Light Amplification by Stimulated Emission of Radiation".
    Con il fatto che i fotoni emessi e quelli che li hanno stimolat sono identici, il laser è una sorgente di luce coerente, ovvero che ha una frequenza, fase e direzione ben definite, e permette di trasmettere dati a velocità elevate su distanze più lunghe.\\
\end{itemize}

Poichè il laser non ha niente in ingresso che lo stimoli, è necessario far in modo che la luce venga riflessa più volte all'interno del laser, in modo da stimolare l'emissione di più fotoni, e quindi creare un fascio di luce coerente.\\
Il laser è costituito da un mezzo attivo, che è un materiale che può essere eccitato per emettere luce, e da una cavità ottica, che è un sistema
di specchi che riflette la luce all'interno del laser, e permette di stimolare l'emissione di più fotoni.\\
Il mezzo attivo può essere costituito da un gas, un liquido o un solido, e viene eccitato da una sorgente di energia, come ad esempio una scarica elettrica, una luce intensa o una reazione chimica, 
e permette di generare luce coerente con una frequenza specifica.\\
La cavità ottica è costituita da due specchi, uno dei quali è parzialmente trasparente (per permettere di trasmettere una parte della luce), e permette alla luce di essere riflessa più volte all'interno del laser, 
e di stimolare l'emissione di più fotoni, e di creare un fascio di luce coerente con una frequenza specifica.\\

\paragraph{LED e Laser a semiconduttore}

Usiamo i Laser a semiconduttore perché costano poco. La corrente elettrica che attraversa il materiale semiconduttore, eccita gli elettroni, rendendo così il mezzo attivo.
I semiconduttori si surriscaldano, e quindi è necessario utilizzare un sistema di raffreddamento per mantenere la temperatura del laser a un livello ottimale, e garantire così prestazioni stabili e affidabili.\\

Nella fiber pigtail esce il laser.\\

Il simbolo del LED è un diodo con due freccettine che indicano che emette luce. Una lacuna si può ricombinare con un elettrone, e quando lo fa emette un fotone, e quindi emette luce.
Un elettrone decade dalla banda di conduzione alla banda di valenza, e quando lo fa emette un fotone, e quindi emette luce.\\
La caratteristica di questo fotone è che l'energia del fotone è uguale alla differenza di energia tra la banda di conduzione e la banda di valenza, 
e quindi la lunghezza d'onda della luce emessa dipende dal materiale semiconduttore utilizzato, e può essere calcolata utilizzando la formula $\lambda = \frac{hc}{E}$, 
dove h è la costante di Planck, c è la velocità della luce e E è l'energia del fotone.\\

Si può approfondire la storia del led blu.\\

Nei LED l'emissione spontanea è dominante, e quindi la luce emessa è incoerente, mentre nei laser a semiconduttore l'emissione stimolata è dominante, 
e quindi la luce emessa è coerente, e permette di trasmettere dati a velocità elevate su distanze più lunghe.\\
Lo spettro di emissione dei LED è più ampio rispetto a quello dei laser a semiconduttore, e quindi i LED emettono luce con una gamma più ampia di lunghezze d'onda,
mentre i laser a semiconduttore emettono luce con una lunghezza d'onda specifica, e quindi sono più adatti per le comunicazioni in fibra ottica.\\

Il diodo laser ha le facce pulite in modo che agiscono come specchi, innietto dall'esterno cariche (quindi elettroni dalla regione P e neutroni dalla regione N) e quindi si crea una popolazione inversa, 
e quindi si innesca l'emissione stimolata, e quindi si genera un fascio di luce coerente con una frequenza specifica.\\

\paragraph{Tipi di Laser}
 %vedere come funzionano
\begin{itemize}
    \item Laser Fabry Perot: è un tipo di laser a semiconduttore che utilizza una cavità ottica costituita da due specchi, uno dei quali è parzialmente trasparente, 
    e permette di generare luce coerente con una frequenza specifica. Sono lo standard per le comunicazioni in fibra ottica a bassa velocità (fino a 1 Gbps) su distanze fino a 2 km.\\
    \item Vertical Cavity Surface Emitting Laser (VCSEL): è un tipo di laser a semiconduttore che utilizza una cavità ottica verticale, e permette di generare luce coerente con una frequenza 
    specifica, e viene utilizzato principalmente nelle comunicazioni in fibra ottica a breve distanza (fino a 300 metri) all'interno dei data center.\\
    \item Distributed Feedback Laser (DFB): è un tipo di laser a semiconduttore che utilizza una cavità ottica costituita da un reticolo di diffrazione, e permette di generare luce coerente con una frequenza specifica,
    e viene utilizzato principalmente nelle comunicazioni in fibra ottica a lunga distanza (oltre 2 km) e ad alta velocità (fino a 100 Gbps), quindi long-haul, metro ...\\
    \item Distributed Bragg Reflector Laser (DBR): è un tipo di laser a semiconduttore che utilizza una cavità ottica costituita da un reticolo di diffrazione, e permette di generare luce coerente con una frequenza specifica,
    e viene utilizzato principalmente nelle comunicazioni in fibra ottica a lunga distanza (oltre 2 km) e ad alta velocità (fino a 100 Gbps).\\
\end{itemize}

Gli DFB e DBR permettono di diminuire la dispersione cromatica e l'ISI.

\paragraph{Modulazione diretta del Laser}

Direct Modulater Laser - DML è una tecnica di modulazione in cui la corrente elettrica che attraversa il laser viene modulata direttamente per generare impulsi di luce, 
e viene utilizzata principalmente nelle comunicazioni in fibra ottica a bassa velocità (fino a 1 Gbps) su distanze fino a 2 km.\\

Il segnale di modulazione viene applicato direttamente a un laser tramite un driver che converte il flusso di dati (tensione) nella corrente di iniezione.

Prima emissione spontanea, quando inizio a mettere più corrente rispetto alla corrente di soglia, inizio ad avere emissione stimolata, e quindi inizio a generare luce coerente 
con una frequenza specifica, e quindi posso trasmettere dati a velocità elevate su distanze più lunghe.\\

\paragraph{Modulazione esterna}
External Modulator Laser - EML è una tecnica di modulazione in cui un modulatore esterno viene utilizzato per modulare la luce generata da un laser, e viene utilizzata principalmente nelle comunicazioni in fibra ottica a lunga distanza (oltre 2 km) e ad alta velocità (fino a 100 Gbps).\\

\paragraph{Modulatori On Off Keying - OOK}

Il vantaggio di questa modulazione è che prendo il segnale così come è e lo mando direttamente al diodo, l'informazione è codificata direttamente nella luce, e quindi è una tecnica di modulazione semplice e efficiente, e viene utilizzata principalmente nelle comunicazioni in fibra ottica a bassa velocità (fino a 1 Gbps) su distanze fino a 2 km.\\
Il vantaggio è che è semplice e efficiente, ma il problema è che non posso raggiungere velocità di trasmissione elevate su distanze più lunghe, perché la modulazione diretta del laser introduce rumore e distorsione nel segnale, e quindi limita la qualità del segnale e la velocità di trasmissione.\\

\paragraph{Modulatori coerenti}

\textbf{Mach-Zehnder Modulator - MZM}\\
Per modulare l'ampiezza e la fase di una  portante si usa un modulatore IQ. In questo caso due segnali in banda base modulano in ampiezza due portanti ortogonali (90 gradi di differenza di fase), e poi le sommo insieme, e quindi ottengo un segnale modulato in ampiezza e fase, e posso raggiungere velocità di trasmissione elevate su distanze più lunghe, perché la modulazione coerente permette di ridurre il rumore e la distorsione nel segnale, e quindi migliorare la qualità del segnale e la velocità di trasmissione.\\

\paragraph{Sistema di trasmissione e ricezione coerenti}

\subsection{Fotodiodi}
Il fotodiodo è un dispositivo che converte la luce in corrente elettrica, e viene utilizzato principalmente nelle comunicazioni in fibra ottica per rilevare la luce trasmessa attraverso la fibra ottica, e convertire il segnale ottico in segnale elettrico, e permette di recuperare le informazioni trasmesse.\\
Usa l'effetto fotoelettrico, che è il fenomeno per cui quando un fotone colpisce un materiale semiconduttore, può liberare un elettrone, e quindi generare una corrente elettrica, 
e permette di convertire la luce in corrente elettrica.\\

\paragraph{Fotodiodo PIN}
Il fotodiodo PIN è un tipo di fotodiodo che utilizza una struttura a strati, costituita da una regione di tipo P, una regione di tipo N e una regione intrinseca (I) tra di esse,
e permette di convertire la luce in corrente elettrica con una maggiore efficienza rispetto ai fotodiodi tradizionali, e viene utilizzato principalmente nelle comunicazioni in fibra
ottica a bassa velocità (fino a 1 Gbps) su distanze fino a 2 km.\\

\paragraph{Fotodiodo Avalanche - APD}
Il fotodiodo Avalanche (APD) è un tipo di fotodiodo che mette in evidenza l'effetto valanga, che è un fenomeno per cui quando un fotone colpisce un materiale semiconduttore, 
può liberare un elettrone, e questo elettrone può a sua volta liberare altri elettroni, creando così una valanga di elettroni, e permette di convertire la luce in corrente elettrica 
con una maggiore sensibilità rispetto ai fotodiodi tradizionali, e viene utilizzato principalmente nelle comunicazioni in fibra ottica a lunga distanza (oltre 2 km) e ad alta 
velocità (fino a 100 Gbps).\\

\paragraph{Materiali per i fotodiodi}
I fotodiodi possono essere realizzati con diversi materiali semiconduttori, come ad esempio il silicio (Si), il germanio (Ge) e l'arseniuro di gallio (GaAs), 
e la scelta del materiale dipende dalle caratteristiche desiderate del fotodiodo, come ad esempio la sensibilità, la velocità di risposta e la lunghezza d'onda della luce da rilevare.\\

\subsection{Amplificatori ottici}
\paragraph{Erbium Doped Fiber Amplifier - EDFA}
Nei collegamenti ottici spesso devo amplificare il segnale ottico, lo faccio quando devo fare le dorsali o submarini perché cerco di mantenere il dato nel dominio ottico il piú possibile.
La categoria di amplificatori ottici più utilizzata è l'Erbium Doped Fiber Amplifier (EDFA), che utilizza una fibra ottica dopata con erbio (terra rara) come mezzo attivo, 
e permette di amplificare il segnale ottico con una bassa attenuazione e un'alta efficienza.
Una fibra ottica viene drogata con ioni Erbio ($Er^{3+}$) per amplificare la radiazione alle lunghezze d'onda in terza finestra ($\lambda = 1550 nm$).
Prende energia dal laser di pompa e amplifica il segnale ottico che viaggia all'interno della fibra ottica, e permette di trasmettere dati a velocità elevate su distanze più 
lunghe senza la necessità di convertire il segnale in elettrico.\\

\paragraph{Amplificatori Raman}
Gli amplificatori Raman utilizzano l'effetto Raman per amplificare il segnale ottico, e consistono in una fibra ottica in cui viene immessa una radiazione laser di pompa, 
che interagisce con il segnale ottico da amplificare, e permette di amplificare il segnale ottico con una bassa attenuazione e un'alta efficienza.\\

\paragraph{Semiconductor Optical Amplifier - SOA}
Il Semiconductor Optical Amplifier (SOA) è un tipo di amplificatore ottico che utilizza un materiale semiconduttore come mezzo attivo, 
e permette di amplificare il segnale ottico con una bassa attenuazione e un'alta efficienza.
Il SOA ha una struttura simile a quella di un diodo laser, ma con trattamenti antiriflesso sulle facce per evitare l'auto-oscillazione, 
e permette di amplificare il segnale ottico con una bassa attenuazione e un'alta efficienza.\\

\paragraph{Planar lightwave circuits - PLC e Photonic Integrated Circuits - PIC}
I Planar Lightwave Circuits (PLC) e i Photonic Integrated Circuits (PIC) sono dispositivi che integrano più componenti ottici su un unico substrato,
e permettono di amplificare il segnale ottico con una bassa attenuazione e un'alta efficienza.
\subsection{Switch ottici}

\paragraph{Switch Micro Electro Mechanical System -MEMS}

\paragraph{Thermo Optic - TO}
\paragraph{Eletro Optic - EO}
\paragraph{SOA}

\paragraph{Optical Cross Connect - OXC}

\section{Wireless}
\subsection{Optical Wireless Communications - OWC}
Sono sistemi di comunicazione che utilizzano la luce per trasmettere dati attraverso l'aria, e possono essere utilizzati in ambienti in cui le comunicazioni in fibra ottica non sono pratiche, come ad esempio in ambienti urbani densi o in ambienti interni.\\

\paragraph{Visible Light Communication - VLC}
La Visible Light Communication (VLC) è una tecnologia di comunicazione wireless che utilizza la luce visibile per trasmettere dati, e può essere utilizzata in ambienti interni, come ad esempio in uffici, case o negozi, per fornire connettività wireless ad alta velocità.\\
Usiamo la luce del visibile perché prendiamo le lampadine a led e le accendiamo e spegniam o con tempi dell'ordine dei nani secondi (all'occhio umano sembra sempre accesa) ma in realtà stiamo trasmettendo un segnale che possiamo usare al posto del wi fi.\\
Sono usati per esempio nei LED di un aereo. 

Il LED bianco non esiste, è l'unione di LED blu rosso e verde che danno l'effetto del colore bianco.
Ha lunghezze d'onda da 390 nm a 750 nm.
\paragraph{Free Space Optical Communication - FSO}
La Free Space Optical Communication (FSO) è una tecnologia di comunicazione wireless che utilizza la luce per trasmettere dati attraverso l'aria, e può essere utilizzata in ambienti esterni, come ad esempio tra edifici o tra città, per fornire connettività wireless ad alta velocità.\\
Comunicazioni tra laser (fotodiodi) montati su edifici, quindi in caso di disastri è possibile ripristinare i collegamenti a volo in free space optical (è più difficile ricostruire l'impianto sotto terra che montare un laser su un edificio).\\
Ha lunghezze d'onda da 750 nm a 1600 nm.
\paragraph{Ultraviolet Communications - UVC}
L'Ultraviolet Communications (UVC) è una tecnologia di comunicazione wireless che utilizza la luce ultravioletta per trasmettere dati, e può essere utilizzata in ambienti esterni, come ad esempio tra edifici o tra città, per fornire connettività wireless ad alta velocità.\\
Viene usato anche quellpo infrarosso. 
Ha lunghezze d'onda da 200 nm a 280 nm.

\paragraph{Light Fidelity - LiFi}
Il Light Fidelity (LiFi) è una tecnologia di comunicazione wireless che utilizza la luce per trasmettere dati, e può essere utilizzata in ambienti interni, come ad esempio in uffici, case o negozi, per fornire connettività wireless ad alta velocità.\\
E' una tecnologia che utilizza la luce visibile per trasmettere dati, e può essere utilizzata in ambienti interni, come ad esempio in uffici, case o negozi, per fornire connettività wireless ad alta velocità.\\
\paragraph{Infrared Communication - IRC}
L'Infrared Communication (IRC) è una tecnologia di comunicazione wireless che utilizza la luce infrarossa per trasmettere dati, e può essere utilizzata in ambienti interni, come ad esempio in uffici, case o negozi, per fornire connettività wireless ad alta velocità.\\
Viene usato anche quello ultravioleto.
Ha lunghezze d'onda da 750 nm a 1600 nm.

\paragraph{Mezzi}
Il VLC usa LED o comunque sorgenti di luce visibile, e fotodiodi come rivelatori di luce, e utilizza protocolli di trasmissione come ad esempio Ethernet o InfiniBand per gestire la comunicazione tra i dispositivi.\\

In OCC ho un insieme di LED o addirittura dei Display, ogni parte del display può inviare informazioni e posso riceverli con una o più camere. Quindi su uno schermo prendo una parte dello schermo e la accendo e spengo con tempi dell'ordine dei nanosecondi, e quindi posso trasmettere un segnale che posso ricevere con una o più camere, e quindi posso trasmettere dati a velocità elevate su distanze più lunghe.\\

MIMO (Multiple Input Multiple Output) è una tecnica di comunicazione wireless che utilizza più antenne per trasmettere e ricevere dati, e può essere utilizzata per migliorare la qualità del segnale e aumentare la velocità di trasmissione nelle comunicazioni in fibra ottica.\\

IRC ho un LED che trasmette luce infrarossa, e un fotodiodo che riceve la luce infrarossa, e utilizza protocolli di trasmissione come ad esempio Ethernet o InfiniBand per gestire la comunicazione tra i dispositivi.\\

IN FSO ho un laser che trasmette luce, e un fotodiodo che riceve la luce, e utilizza protocolli di trasmissione come ad esempio Ethernet o InfiniBand per gestire la comunicazione tra i dispositivi.\\

In UVC ho un LED che trasmette luce ultravioletta, e un fotodiodo che riceve la luce ultravioletta, e utilizza protocolli di trasmissione come ad esempio Ethernet o InfiniBand per gestire la comunicazione tra i dispositivi.\\

Se io mando il LASER nello spazio libero dovrei preoccuparmi della sicurezza. L'occhio è sensibile alla luce visibile, ma se lavoriamo nell'infrared posso dare danni all'occhio senza che la persona se ne accorga. Se focalizzo la luce nell'infrarosso creo danni permanenti, posso rischiare di cuocere la cornea.

Usiamo la luce perché la luce non è normata (siamo liberi di usare qualunque frequenza d'onda senza dover pagare le licenze), e perché la luce ha una banda molto ampia, e quindi permette di trasmettere dati a velocità elevate su distanze più lunghe. Per questo la preferiamo alle radiofrequenze, che sono normate e hanno una banda limitata, e quindi non permettono di trasmettere dati a velocità elevate su distanze più lunghe.\\

Multipath Fading è un fenomeno che si verifica quando un segnale wireless viene riflesso da oggetti come edifici, alberi o veicoli, e quindi arriva al ricevitore con più percorsi diversi, e può causare interferenze e distorsioni nel segnale, e quindi limitare la qualità del segnale e la velocità di trasmissione nelle comunicazioni wirelesss. È un problema delle radiofrequenze e non della luce, perché la luce non viene riflessa dagli oggetti, e quindi non subisce il fenomeno del multipath fading.\\

\paragraph{Standard di sicurezza dei laser}
I laser sono classificati in base alla loro potenza e al rischio che rappresentano per la salute umana, e sono suddivisi in 4 classi principali:
\begin{itemize}
    \item Classe 1: sono laser a bassa potenza che non rappresentano un rischio per la salute umana, e possono essere utilizzati senza restrizioni.\\
    \item Classe 2: sono laser a bassa potenza che rappresentano un rischio per la salute umana se vengono guardati direttamente per un periodo prolungato, e devono essere utilizzati con precauzioni.\\
    \item Classe 2M: sono laser a bassa potenza che rappresentano un rischio per la salute umana se vengono guardati direttamente con strumenti ottici, come ad esempio binocoli o telescopi, e devono essere utilizzati con precauzioni.\\
    \item Classe 3: sono laser a media potenza che rappresentano un rischio per la salute umana se vengono guardati direttamente, e devono essere utilizzati con precauzioni e dispositivi di protezione.\\
    \item Classe 3B: sono laser a media potenza che rappresentano un rischio per la salute umana se vengono guardati direttamente, e devono essere utilizzati con precauzioni e dispositivi di protezione, e sono soggetti a restrizioni legali.\\
    \item Classe 3R: sono laser a media potenza che rappresentano un rischio per la salute umana se vengono guardati direttamente, e devono essere utilizzati con precauzioni e dispositivi di protezione, e sono soggetti a restrizioni legali.\\
    \item Classe 4: sono laser ad alta potenza che rappresentano un rischio per la salute umana se vengono guardati direttamente, e devono essere utilizzati con precauzioni e dispositivi di protezione, e sono soggetti a restrizioni legali.\\
\end{itemize}

\paragraph{OWC vs Wireless (RF)}

\paragraph{Caratteristiche dei sistemi OWC}

Il processamento del segnale elettrico è in banda base, trasformo il segnale da elettrico in ottico. Quando ricevo metto un fotodiodo nel field of view del laser. Questo field of view è più stretto per i laser perché sono più focalizzati rispetto alla luce LED. Quindi il fotodiodo sarà più piccolo per i laser rispetto ai LED quindi lavoro Line of Sight (LOS) con i laser, e non Line of Sight (NLOS) con i LED.\\
Il vantaggio di lavorare in LOS è che posso raggiungere velocità di trasmissione elevate su distanze più lunghe, perché il segnale ottico non subisce interferenze e distorsioni causate da oggetti come edifici, alberi o veicoli, e quindi posso trasmettere dati a velocità elevate su distanze più lunghe.\\
Il vantaggio di lavorare in NLOS è che posso trasmettere dati anche in ambienti in cui la linea di vista è ostruita da oggetti come edifici, alberi o veicoli, e quindi posso fornire connettività wireless in ambienti urbani densi o in ambienti interni, ma il problema è che il segnale ottico subisce interferenze e distorsioni causate da oggetti come edifici, alberi o veicoli, e quindi limita la qualità del segnale e la velocità di trasmissione nelle comunicazioni wireless.\\

I ricevitori posso usarli di qualunque tipo. 

Il vantaggio di OWC è che ho una banda molto grande, è sicuro perché la luce non attraversa i muri quindi le informazioni rimangono nell'ambiente, può essere utilizzato negli ambieniti come ospedali aerei o impianti industriali perché non c'è rischio di interferenze elettromagnetiche, ed è a basso costo.\\
Gli svantaggio sono il Los quindi gli ostacoli impediscono la comunicazione, L'interferenza atmosferica come luce solare o artificiale che introducono rumore, l'influenza meterologica (FSO), e una mobilità ridotta.\\

\paragraph{Applicazioni dei sistemi OWC}
Usati negli ospedali, negli aerei e conferenze, nelle industrie, nelle reti V2V e trasporto, nel backhaul, nelle comunicazioni militari (FSO e UV).\\

\subsection{Optical wireless PAN e BAN}

Personal Area Network (PAN) è una rete di comunicazione che collega dispositivi personali, come ad esempio smartphone, tablet, laptop e dispositivi indossabili, all'interno di un'area limitata, come ad esempio una stanza o un edificio.\\
Body Area Network (BAN) è una rete di comunicazione che collega dispositivi indossabili o impiantati all'interno del corpo umano, come ad esempio sensori medici, dispositivi di monitoraggio della salute e dispositivi di assistenza sanitaria, per fornire connettività wireless e supportare applicazioni come la telemedicina e la salute digitale.\\
\subsection{Optical wireless Vehicular Ad Hoc Network - VANET}
Sono reti che comunicano tra due veicoli o tra veicoli e infrastrutture stradali, e possono essere utilizzate per fornire connettività wireless ad alta velocità e supportare applicazioni come la guida autonoma, la sicurezza stradale e la gestione del traffico.\\

\subsection{ LANET}

\subsection{Optical Camera Communications - OCC}

\paragraph{OWC per applicazioni aeronautiche}

\paragraph{Standards}
PPM (Pulse Position Modulation) è una tecnica di modulazione in cui la posizione di un impulso all'interno di un intervallo di tempo viene utilizzata per rappresentare un simbolo. Quindi in base alla posizione dell'impulso (che trasmetto sempre) so se è uno 0 o 1.\\

ciao

Tecniche di modulazione per OWC:
\begin{itemize}
    \item OOK
    \item OFDM
    \item PPM
    \item PAM-4
\end{itemize}

In genere voglio o ricevere la luce LoS o NLoS, se uso tutte e due posso imbattermi in delle interferenze.

Confronto LiFi vs WiFi:

\subsection{Free Space Optics - FSO}

\subsection{Underwater optical wireless communications}
L'acqua attenua di meno la luce blu e molto la luce rossa, quindi per le comunicazioni sottomarine si usano laser blu o verdi.\\

Nell'acqua la velcità della luce è circa 200 mila km/s, quindi è più lenta che nell'aria, e questo può causare problemi di sincronizzazione e di latenza nelle comunicazioni sottomarine.\\

Il problema grosso di queste comunicazioni è l'attenuazione.




\chapter{Reti in fibra ottica}
\section{Pulse Code Modulation (PCM)}
Il segnale analogico viene campionato a una frequenza di campionamento che deve essere almeno il doppio della frequenza massima del segnale (teorema di Nyquist), e quindi quantizzato in un certo numero di livelli, e quindi codificato in un certo numero di bit per ogni campione.\\
Il PCM è un metodo di codifica digitale che viene utilizzato per rappresentare i segnali analogici in forma digitale, e quindi permette di trasmettere i segnali su reti digitali, e quindi migliorare la qualità del segnale e aumentare la velocità di trasmissione, e supporta anche la trasmissione bidirezionale, e quindi permette di fornire connettività wireless ad alta velocità in ambienti in cui le comunicazioni in fibra ottica non sono pratiche.\\

Per il segnale voce la banda è a 4 kHz, quindi per il teorema di Nyquist la frequenza di campionamento deve essere almeno 8 kHz, e quindi con una quantizzazione a 8 bit per campione si ottiene una velocità di trasmissione di 64 kbps, si trasmette un campione ogni 125 µs (1/8000).\\

Quantizzazione - Ogni campione viene rappresentato da un numero binario di 8 bit, e quindi si ottiene una rappresentazione digitale del segnale analogico, e quindi permette di trasmettere i segnali su reti digitali, e quindi migliorare la qualità del segnale e aumentare la velocità di trasmissione, e supporta anche la trasmissione bidirezionale, e quindi permette di fornire connettività wireless ad alta velocità in ambienti in cui le comunicazioni in fibra ottica non sono pratiche.\\

Il segnale digitale viene interllacciato con i segnali provenienti da altri utenti per formare una Digital Hierarchy. Questa tecnica si chiama TDM (Time Division Multiplexing).\\

\section{Plesiochronous Digital Hierarchy (PDH)}
Il PDH è una tecnica di multiplexing che consente di combinare più segnali digitali a velocità diverse in un unico flusso di dati, e quindi permette di trasmettere i segnali su reti digitali, e quindi migliorare la qualità del segnale e aumentare la velocità di trasmissione, e supporta anche la trasmissione bidirezionale, e quindi permette di fornire connettività wireless ad alta velocità in ambienti in cui le comunicazioni in fibra ottica non sono pratiche.\\

Plesiocrane significa che il clock non era lo stesso. 

\section{Synchronous Digital Hierarchy (SDH) - Synchronous Optical Network (SONET)}
Fu standardizzata intorno agli '90, in italia dalla telecom. Siamo passati dalle reti PDH a quelle SDH/SONET per aggiungere controllo delle rete. LA rete è costituita da una serie di trasmettitori che li mette insieme dal multiplexing per creare il segnale di trasporto. Lungo la rete abbiamo degli amplificatori. ADM è un nodo.
STS DEMUX separa poi il segnale.\\

Questi 3 livelli di overhead mi permette un controllo del flusso, mi permette di aumentare la velocità in fibra. TUtti gli elementi della rete usano lo stesso clock, sono sincronizzati in ampiezza e in fase. Avere questa struttura ben organizzata ci permette di aprire il mio switch a determinati istanti di tempo ed estrarre il flusso di dati che mi interessa.\\

Ho informazioni sullo stato della rete, posso fare il monitoraggio, posso fare il controllo del flusso, posso fare la gestione degli errori.\\

Introduco quello che si chiama la qualità del servizio. Quindi l'operatore oltre a garantirti la banda, ti garantisce anche la qualità di servizio.\\

Abbiamo una struttura stratificata. Prendiamo il nostro segnale elettrico, viene inserito il mio dato e aggiungo un overhead di cammino, dove questo overhead viene letto solo ricezione. Quando faccio il multiplexing ci metto l'overhead di linea, quindi mi controlla la sincronizzazione e l'errore nella comunicazione. Poi aggiungo overhead di sezione, e infine trasformo il segnale elettrico in ottico (photonic).\\

Tre tipi di overhead:
\begin{itemize}
    \item Di cammino
    \item Di linea
    \item Di sessione
\end{itemize}

\paragraph{SONET: Frame STS-n}


\chapter{IoT}
\section{Internet of things}
Ci occupiamo dei sistemi di connessione compresi bluethoth, wifi e LoRaWAN \dots\\
Quello che studiamo sono, non le applicazioni dei sistemi IoT o dei sensori stessi, ci occupiamo della trasmissione dei dati.
Il bluethooth detto classic (BR per Basic Rate, EDR per Enhanced Data Rate) è un sistema di connessione a corto raggio, con una portata massima di 100 metri, e una velocità massima di 3 Mbps. La versione 4.0 è definita smart bluethooth o low energy (BLE), la BT (bluethooth classic) e BLE hanno protocolli diversi quindi non sono compatibili, i nostri dispositivi li usano entrambi.\\

Le long Range Wide Area Network (LoRaWAN), le reti a long range (dimensioni a qualche chilometro) stanno proliferanno, stanno diventando più comuni e quindi è probabile che li vedremo nella nostra vita lavorativa.\\

Usiamo delle tecnologie esistenti o estensioni delle tecnologie esistenti, le dividiamo inoltre in base all'uso.

La telefonia viene utilizzata per altre applicazioni, tipicamente non vengono utilizzate per IoT, usare le frequenze per la telefonia è limitante, per questo per IoT vengono usate frequenze non licenziate chiamate ISM (Industrial, Scientific and Medical) che sono libere da licenze, e quindi possono essere utilizzate da chiunque, e sono più adatte per IoT.\\

SI può modificare una rete LTE per farla diventare IoT, in questo caso si parla di LTE-M (LTE for Machine Time Comunication), è una tecnologia che utilizza le frequenze della telefonia, ma è stata progettata per essere più efficiente in termini di consumo energetico e di costi, rispetto alla LTE tradizionale.\\

Nelle LTE cerchiamo di lavorare nelle frequenze GSM, il segnale si propaga con distanze ampie, e poiché le reti Low power usano lunghe distanze a costo ridotto.\\
Le reti Low Power utilizzano Unlicensed frequenze.
Bluethooth lo usiamo per sostituire il cavo, per connessioni di breve raggio, con una banda decente. 

Zigbee è un protocollo di comunicazione wireless basato su IEEE 802.15.4, è progettato per essere utilizzato in reti a bassa potenza e a corto raggio, è utilizzato principalmente per applicazioni di automazione domestica e industriale.\\

Definiamo oltre agli standard anche le applicazioni. Le personal/ local area network (PAN/LAN) sono reti che coprono un'area limitata, come una casa o un ufficio, e sono utilizzate per connettere dispositivi all'interno di quella area, possono funzionare con poca potenza.\\

Le Low Power Wide Area Network (LPWAN) sono reti a bassa potenza e a lunga distanza, sono progettate per connettere dispositivi che si trovano a distanze maggiori rispetto alle PAN/LAN, e sono utilizzate principalmente per applicazioni di monitoraggio e controllo remoto.\\

Le cellular network sono reti che utilizzano le frequenze della telefonia, sono progettate per connettere dispositivi che si trovano a distanze maggiori rispetto alle PAN/LAN e LPWAN, e sono utilizzate principalmente per applicazioni di monitoraggio e controllo remoto.\\

\subsection{Low Power Wide Area Network - LPWAN}
Se devo misurare per esempio un contatore dell'acqua mi serve la rete, e alcune compagnie stanno usando le LPWAN per questo, è un'applicazione tipica, ma non è l'unica.\\
Ormai è la tecnologia più usata, è una tecnologia che sta crescendo molto, e ci sono molte aziende che stanno investendo in questo settore.\\
Hanno due caratteristiche principali, la prima è la lunga distanza, che può arrivare fino a 10 km in ambiente urbano e fino a 15 km in ambiente rurale, e la seconda è il basso consumo energetico, che permette ai dispositivi di funzionare per anni con una sola batteria.\\

Le LPWAN utilizzano frequenze non licenziate, come la banda ISM a 868 MHz in Europa e a 915 MHz negli Stati Uniti, e sono progettate per essere utilizzate in ambienti urbani e rurali.\\

Le LPWAN sono utilizzate principalmente per applicazioni di monitoraggio e controllo remoto, come il monitoraggio dei contatori dell'acqua, il monitoraggio dei livelli di inquinamento, il monitoraggio dei livelli di rumore, il monitoraggio dei livelli di temperatura, il monitoraggio dei livelli di umidità, il monitoraggio dei livelli di luce, il monitoraggio dei livelli di movimento, il monitoraggio dei livelli di vibrazione, il monitoraggio dei livelli di pressione, il monitoraggio dei livelli di gas, il monitoraggio dei livelli di fumo, il monitoraggio dei livelli di radiazione, il monitoraggio dei livelli di umidità del suolo, il monitoraggio dei livelli di umidità dell'aria, il monitoraggio dei livelli di umidità del terreno, il monitoraggio dei livelli di umidità del pavimento, il monitoraggio dei livelli di umidità del soffitto, il monitoraggio dei livelli di umidità del muro, il monitoraggio dei livelli di umidità del tetto, il monitoraggio dei livelli di umidità del pavimento, il monitoraggio dei livelli di umidità del soffitto, il monitoraggio dei livelli di umidità del muro, il monitoraggio dei livelli di umidità del tetto.\\

Le LPWAN sono utilizzate anche per applicazioni di automazione domestica e industriale, come il controllo delle luci, il controllo delle porte, il controllo delle finestre, il controllo dei termostati, il controllo dei sistemi di sicurezza, il controllo dei sistemi di irrigazione, il controllo dei sistemi di ventilazione, il controllo dei sistemi di riscaldamento, il controllo dei sistemi di raffreddamento, il controllo dei sistemi di illuminazione, il controllo dei sistemi di allarme, il controllo dei sistemi di sorveglianza, il controllo dei sistemi di monitoraggio, il controllo dei sistemi di gestione dell'energia, il controllo dei sistemi di gestione dell'acqua, il controllo dei sistemi di gestione del gas, il controllo dei sistemi di gestione dei rifiuti.\\

\subsection{Personal Area Network - PAN}
Studiamo lo standard Bluetooth, è uno standard di comunicazione wireless a corto raggio, progettato per connettere dispositivi personali come smartphone, tablet, laptop, cuffie, altoparlanti, stampanti, tastiere, mouse, e altri dispositivi.\\
Tutte le reti PAN devono utilizzare bassa potenza, dalla versione 4.0 Bluetooth è diventata Low Enerrgy, la topologia è a stella nelle bluetooth standard, ma è a mesh nelle bluetooth low energy, la banda è di 2.5 GHz, e la velocità massima è di 3 Mbps per Bluetooth standard e di 1 Mbps per Bluetooth low energy.\\

Il wifi e bluetooth lavorano entrambi a 2.5 GHz, ma il wifi ha una banda più ampia e una velocità maggiore, quindi è più adatto per applicazioni che richiedono una maggiore larghezza di banda, come lo streaming video o il trasferimento di file di grandi dimensioni, mentre il bluetooth è più adatto per applicazioni che richiedono una connessione a corto raggio e un basso consumo energetico, come la connessione di dispositivi personali. Quindi poiché il wifi è più dominante occupa dei canali e il bluetooth deve adattarsi, per questo il bluetooth ha una banda più stretta e una velocità inferiore per evitare interferenze con il wifi.\\

Il bluetooth è uno standard SIG (Special Interest Group) che è stato sviluppato da un consorzio di aziende, tra cui Ericsson, IBM, Intel, Nokia, e Toshiba.\\

Si usa perché è molto facile accoppiare i dispositivi, e la connessione è diretta ed affidabile e trasmette a 2.5 GHz, quindi è meno soggetta a interferenze rispetto ad altre frequenze.\\

Lo standard è 802.15, 802.3 è ethernet, 802.11 è wifi.\\

Lo utilizziamo per sostituire i cavi (stampanti, cuffie, altoparlanti, tastiere, mouse), per connettere dispositivi personali, e per applicazioni di automazione domestica e industriale. Come raggio ha brevi distanze (sotto i 10 metri), a basso costo, e a basso consumo energetico. Sono nati per sostituire l'infrared che aveva una portata molto limitata e richiedeva una linea di vista diretta tra i dispositivi.\\
Se ho un flag che posso settare per far si che ti informo che ho ricevuto il pacchetto, se non succede ti rimando il pacchetto in maniera automatica, questo è un meccanismo di affidabilità, e se non riesco a consegnare il pacchetto dopo un certo numero di tentativi, allora ti informo che non sono riuscito a consegnare il pacchetto.\\

\paragraph{Bluetooth}
Il bluetooth 4.0 lo ha creato la nokia. Nel classic la rete è a stella (master con 7 slave, quindi il master si può collegare a 7 dispositivi), mentre nel low energy la rete è a mesh, quindi ogni dispositivo può collegarsi a più dispositivi, e i dispositivi possono comunicare tra loro senza dover passare attraverso un master.\\
Il bluetooth low energy è stato progettato anche per connettere i cartellini elettroni del supermercato, è con low energy riesco a misurare la provenienza del segnale, cosa che non riesco a fare con il bluetooth standard, e questo è utile per applicazioni di localizzazione.\\

Nel bluetooth low energy, cambio completamente la caratteristica dello spettro, infatti nello standard ho banda che suddivido in 79 canali di un MHz, quando trasmetto di nuovo cambio canale, e questo mi permette di evitare interferenze con il wifi, che occupa una banda più ampia, nel BLE ho meno canali (40) ma stessa banda quindi posso essere più flessibile.\\
Uso la codifica per correggere l'errore, quindi diminuisco l'informazione ma aggiungo bit per allungare e arrivare a distanze maggiori, quindi con segnale rumoroso ma che posso correggere.\\
\paragraph{Bluetooth classic}
la rete bluetooth classic si chiama piconet.
Il master può collegarsi  a massimo 7 dispositivi, perché ogni dispositivo viene individuato nella rete da un indirizzo di 3 bit, quindi 2 alla potenza di 3 è uguale a 8, ma uno è riservato per il master, quindi rimangono 7 indirizzi per gli slave.\\

Scatternet è un insieme di piconet, quindi posso avere più piconet che si collegano tra loro, e i dispositivi possono comunicare tra loro senza dover passare attraverso un master.\\

Non trasmetto solo alto o basso (0 o 1), uso la frequency shift keying, fisicamente mando un segnale che oscilla a una frequenza, se voglio trasmettere 0 mando un segnale che oscilla a una frequenza, se voglio trasmettere 1 mando un segnale che oscilla a un'altra frequenza, e questo mi permette di trasmettere informazioni in modo più efficiente e affidabile.\\

I pacchetti durano da 366 microsecondi a 625 microsecondi, e la velocità massima è di 3 Mbps, ma in realtà la velocità effettiva è inferiore a causa dell'overhead dei pacchetti e delle interferenze.\\ 
La durata di ogni pacchetto può essere numero dispari di slot, e ogni slot dura 625 microsecondi, quindi la durata minima di un pacchetto è di 366 microsecondi (1 slot) e la durata massima è di 625 microsecondi (2 slot).\\
Posso trasmettere 1 slot, che dura 3 slot o 5 slot, inoltre il master trasmette sempre nei slot pari, mentre gli slave trasmettono nei slot dispari, questo permette di evitare collisioni tra i pacchetti.\\ 

Non usa una sola frequenza, se quella frequenza si blocca (per esempio c'è il wifi) allora cambio frequenza, e questo mi permette di evitare interferenze e di migliorare l'affidabilità della connessione. Ogni volta che trasmetto un pacchetto, la mia frequenza cambia, c'è salto di frequenza, questo cambio è random. Sto cambiando la frequenza 1600 volte al secondo, quindi ogni 625 microsecondi, e questo mi permette di evitare interferenze con il wifi, che occupa una banda più ampia.\\
Se i pacchetti hanno durata di 3 slot, la frequenza viene variata di 1600/3 secondi.\\

Quindi il master e lo slave si devono sincronizzare, che è necessario affinché si aggancino nello stesso clock, poi devo decidere quale dei 79 canali usare (lo decide il master) con la tecnica del salto di frequenza, e poi devo decidere quando trasmettere, quindi devo sincronizzare i tempi di trasmissione. Il master decide quale algoritmo randomico mi permette di cambiare frequenza. L'algoritmo di cambio di frequenza di chiama Frequency Hopping Spread Spectrum (FHSS), e questo mi permette di evitare interferenze con il wifi, che occupa una banda più ampia.\\

\paragraph{Frequency Hopping Spread Spectrum - FHSS}
Nel piano di frequenza faccio dei salti pseudorandom, ovvero cambio frequenza con generazione random ma questa generazione è conosciuta sia dal master che dallo slave, quindi è pseudorandom, e questo mi permette di evitare interferenze con il wifi, che occupa una banda più ampia.\\

Stata ideata da Hedy Lamarr, un'attrice austriaca che ha inventato questa tecnica durante la seconda guerra mondiale per evitare che i segnali radio venissero intercettati dai nemici.\\

Nei bluetooth di terza generazione si usa l'\textbf{Adaptive Frequency Hopping (AFH)}, che è una tecnica di salto di frequenza adattiva, che permette di evitare le frequenze che sono occupate da altri dispositivi, come il wifi, e di migliorare l'affidabilità della connessione. Questo range non lo uso perché magari c'è il wifi, quindi faccio salti senza usare i "bad channel" ovvero canali che potrebbero generare interferenze e quindi rumore. Come tutti i dispositivi 802 sono identificati da un indirizzo MAC, che è un identificatore univoco a 48 bit, e questo mi permette di identificare i dispositivi nella rete. Prendo una parte dell'indirizzo MAC e li trasmetto agli slave, in questo modo gli slave possono identificare il master e sincronizzarsi con lui avendo le informazioni su come cambiare la frequenza.\\

Per capire il next Hop, misuro il clock e ne aggiungo un offset, in questo modo posso prevedere quando il master trasmetterà il prossimo pacchetto, e posso sincronizzarmi con lui, quindi ho bisogno innanzitutto che master e slave siano allo stesso clock. Incremento il mio clock di 1 ogni 312,5 microsecondi, e quando ricevo un pacchetto, prendo il clock del master e lo sincronizzo con il mio clock, in questo modo posso prevedere quando il master trasmetterà il prossimo pacchetto, e posso sincronizzarmi con lui. Il master e lo slave trasmettono a vicenda il proprio clock in modo che l'altro si aggiusta con il proprio offset. 

Tutti i dispositivi 802 hanno un \textbf{Media Access Control - MAC} è costituito da 48 bit, i 48 bit di solito vengono rappresentati in esadecimale, quindi ogni esadecimale rappresenta 4 bit, ma poichè è costituito da 48 bit, viene rappresentato da 12 esadecimali, e questo mi permette di identificare i dispositivi nella rete.\\
L'indirizzo MAC è costituito da due parti, Vendor ID e Vendor Allocated. Il Vendor ID è costituito dai primi 24 bit, e identifica il produttore del dispositivo, mentre il Vendor Allocated è costituito dagli ultimi 24 bit, e viene assegnato dal produttore in modo univoco a ogni dispositivo.\\
Il MAC viene diviso in 3 parti, LAP, UAP e NAP. Il LAP (Lower Address Part) è costituito dai primi 24 bit, e viene utilizzato per identificare il dispositivo nella rete, l'UAP (Upper Address Part) è costituito dai successivi 8 bit, e viene utilizzato per identificare il dispositivo nella rete, e il NAP (Non-significant Address Part) è costituito dagli ultimi 16 bit, e viene utilizzato per identificare il dispositivo nella rete.\\
Il MAC del master (UAP) viene utilizzato per definire l'algoritmo che mi genera in modo pseudorandom l'hop. Per definire quale piconet sto definisco il LAP del master, e questo mi permette di identificare la piconet, e di sincronizzarmi con il master.\\

\textbf{Connessioni}
Asincrone o sincrone. Nessuno slave si azzarda a trasmettere se non ha ricevuto un pacchetto dal master, quindi è necessario riservare banda per una connessione sincrona. Il master trasmette e crea una connessione sincrona con lo slave 1 (uno slave può avere solo una connessione sincrona).\\
Se invece uso una connessione asincrona posso creare più connessioni verso slave diverse, invede per le sincrone devono essere bidirezionali. 
Nelle connessioni sincrone, poichè la velocità deve essere 64 bit/s, non si usa nè FEC nè ARQ che invece sono usate nelle asincrone.\\

\textbf{Pacchetti}
I pacchetti hanno una durata di 5 slot, se prendo 5 e lo moltiplico per 625 microsecondi, ottengo 3.125 millisecondi. È sempre costituito da un Access Code, un Header, un Payload. 
L'Access Code è costituito da 72 bit, e viene utilizzato per identificare la piconet, e per sincronizzarsi con il master. L'Header è costituito da 54 bit, e viene utilizzato per identificare il tipo di pacchetto, e per identificare il dispositivo di destinazione. Il Payload è costituito da un numero variabile di bit, e viene utilizzato per trasmettere i dati.\\
L'Access code è sempre trasmessa con un basic rate perché altrimenti non saprei come leggerli, e questo mi permette di identificare la piconet, e di sincronizzarmi con il master. L'Header è sempre trasmesso con un basic rate perché altrimenti non saprei come leggerli, e questo mi permette di identificare il tipo di pacchetto, e di identificare il dispositivo di destinazione. Il Payload può essere trasmesso con un enhanced data rate o con un basic rate, a seconda del tipo di pacchetto, e questo mi permette di trasmettere i dati in modo più efficiente e affidabile.\\



\paragraph{Bluetooth low energy}

\chapter{Sistemi Mobili}
\section{Reti di telefonia mobile}
All'inizio, quindi anni 90, abbiamo ogni 10 anni una nuova generazione di reti mobili, con un salto tecnologico importante. Oggi siamo alla 5G, ma è già in fase di sviluppo la 6G.

TACS (Total Access Communication System) è stata la prima generazione di reti mobili, introdotta negli anni 80. Era una rete analogica, che permetteva solo chiamate vocali e aveva una capacità limitata.
AMPS (Advanced Mobile Phone System) è stata la prima rete mobile digitale, introdotta negli anni 80. Permetteva chiamate vocali e messaggi di testo, ma aveva ancora una capacità limitata.\\

Negli anni 90 sono stati introdotti i primi standard di rete mobile digitale, come GSM (Global System for Mobile Communications) con cui abbiamo avuto connessioni digitali, avevamo sia timeslot che frequenza.
EDGE (Enhanced Data rates for GSM Evolution) è stato un miglioramento del GSM, introdotto negli anni 2000. Permetteva velocità di trasmissione dati più elevate, fino a 384 kbps, ma era ancora limitato rispetto alle reti successive.
GPRS (General Packet Radio Service) è stato un miglioramento del GSM. Permetteva la trasmissione di dati a pacchetto, con velocità fino a 114 kbps, ma era ancora limitato rispetto alle reti successive.

Quando abbiamo introdotto i sistemi GPRS a EDGE siamo passati dalla seconda generazione alla 2.5.\\
La terza generazione fa riferiferimento a UMTS (Universal Mobile Telecommunications System), introdotto negli anni 2000. Permetteva velocità di trasmissione dati più elevate, fino a 2 Mbps, e supportava servizi multimediali come videochiamate e streaming.
Si assegnano codici e non piú timeslot o frequenze.
HSPA (High Speed Packet Access) è stato un miglioramento del UMTS, introdotto negli anni 2000. Permetteva velocità di trasmissione dati più elevate, fino a 14 Mbps, ma era ancora limitato rispetto alle reti successive.
Si passa da una tecnica ibrida time division multiple access (TDMA) e code division multiple access (CDMA) a una tecnica completamente basata su CDMA.\\
La quarta generazione fa riferimento a LTE (Long Term Evolution), introdotto negli anni 2010. Permetteva velocità di trasmissione dati più elevate, fino a 100 Mbps, e supportava servizi multimediali avanzati come video in alta definizione e giochi online.
LTE non é i sistemi 4G, per dire che é un sistema 4G allora stiamo parlando di LTE-Advanced, che permette velocità di trasmissione dati fino a 1 Gbps.\\
La quinta generazione fa riferimento a 5G, introdotto negli anni 2020. Permette velocità di trasmissione dati molto elevate, fino a 10 Gbps, e supporta servizi multimediali avanzati come realtà virtuale e aumentata, e l'Internet of Things (IoT).
Il 5G si basa su 3 pilastri, l'EMBB (Enhanced Mobile Broadband) per la connettività ad alta velocità, l'URLLC (Ultra-Reliable Low Latency Communications) per applicazioni che richiedono bassa latenza e alta affidabilità, e mMTC (massive Machine Type Communications) per supportare un gran numero di dispositivi connessi, come quelli dell'IoT.
Il 5G non ha un vero nome ma lo si chiama New Radio (NR).\\

I sistemi 6G non sono ancora stati standardizzati, ma si prevede che offriranno velocità di trasmissione dati ancora più elevate, fino a 100 Gbps, e supportino servizi multimediali avanzati come realtà virtuale e aumentata, e l'Internet of Things (IoT) in modo ancora più efficiente. Si prevede che i sistemi 6G utilizzeranno tecnologie come l'intelligenza artificiale e la comunicazione quantistica per migliorare ulteriormente le prestazioni delle reti mobili.

\subsection{1G}
Martin Cooper inventore del primo telefono cellulare, ha effettuato la prima chiamata da un telefono cellulare il 3 aprile 1973. Il telefono utilizzato era un Motorola DynaTAC, che pesava circa 1 kg e aveva un'autonomia di circa 30 minuti in conversazione.
In realtá i primi cellulari erano disponibili per le auto che si collegavano alla batteria, infatti il problema principale era creare un telefono con autonomia di batteria.
Si chiama cellulare perché la copertura della rete è suddivisa in celle, ognuna delle quali è servita da una stazione base. Quando un utente si sposta da una cella all'altra, la chiamata viene trasferita automaticamente alla nuova stazione base, garantendo una continuità di servizio.
Le celle avevano una dimensione fissa di circa 10 km, e la capacità di ogni cella era limitata a circa 100 utenti. Questo ha portato a problemi di congestione nelle aree urbane, dove la domanda di servizi mobili era elevata.\\
La posizione dell'utente all'interno delle celle era gestita da un sistema piú o meno automatico, che utilizzava segnali radio per determinare la posizione dell'utente e instradare le chiamate in modo efficiente.
L'assegnazione delle varie frequenze era gestita da un sistema di controllo centrale, che assegnava le frequenze alle stazioni base in base alla domanda e alla disponibilità. Questo sistema era basato su un algoritmo di assegnazione delle frequenze che cercava di minimizzare l'interferenza tra le celle adiacenti. Quando io mi sposto da una cella all'altra, il sistema di controllo centrale gestisce il processo di handover, che consiste nel trasferire la chiamata dalla stazione base della cella precedente a quella della nuova cella. Questo processo è gestito in modo da garantire una continuità di servizio e minimizzare l'interruzione della chiamata.\\
Quando questo sistemo é stato sviluppato negli anni 80, l'apparecchio e fare o ricevere le telefonate era molto costoso, per questo si limitavano le chiamate a quelle più importanti, e si preferiva utilizzare i telefoni fissi per le chiamate quotidiane. Inoltre, la copertura della rete era limitata, soprattutto nelle aree rurali, il che ha ulteriormente limitato l'adozione dei telefoni cellulari.
Il sistema é morto perché lo hanno clonato, quando l'utente trasmetteva e si agganciava trasmetteva in chiaro il suo numero di telefono e sia il numero identificativo del telefono e questo causava problemi di sicurezza, infatti era possibile clonare il telefono e fare chiamate a spese dell'utente originale. 
\subsection{2G}
La seconda generazione di reti mobili, introdotta negli anni 90, ha portato a una rivoluzione nella comunicazione mobile. Con l'introduzione del GSM (Global System for Mobile Communications), le chiamate vocali sono diventate digitali, migliorando la qualità del suono e la sicurezza delle comunicazioni. Inoltre, il GSM ha introdotto la possibilità di inviare messaggi di testo (SMS), che è diventato un modo popolare per comunicare brevemente e rapidamente. La capacità della rete è aumentata significativamente, consentendo a più utenti di connettersi contemporaneamente senza problemi di congestione. Tuttavia, il GSM aveva ancora limitazioni in termini di velocità di trasmissione dati, che sono state superate con l'introduzione di tecnologie come GPRS ed EDGE, che hanno portato alla transizione verso la 2.5G.\\
Hanno inventato la possibilitá di identificare il mobile terminal con la SIM (Subscriber Identity Module), che è una scheda rimovibile che contiene informazioni sull'abbonato, come il numero di telefono e le credenziali di autenticazione. Questo ha permesso agli utenti di cambiare telefono mantenendo lo stesso numero di telefono e ha migliorato la sicurezza delle comunicazioni, poiché le informazioni dell'abbonato erano memorizzate in modo sicuro sulla SIM.
La tecnica di modulazione era Gaussian Minimum Shift Keying (GMSK), che è una tecnica di modulazione digitale che utilizza una forma d'onda gaussiana per modulare il segnale. Questa tecnica è stata scelta per il GSM perché offre una buona efficienza spettrale e una bassa interferenza tra i canali adiacenti, consentendo di utilizzare le frequenze in modo più efficiente e migliorando la qualità delle comunicazioni.\\
La trasmissione era digitale e usava insieme sia frequency division multiple access (FDMA) che time division multiple access (TDMA). FDMA divide lo spettro di frequenza in canali separati, mentre TDMA divide il tempo in slot, consentendo a più utenti di condividere la stessa frequenza. Se uso TDMA significa che devo avere un aggancio sincronizzato, se invece uso FDMA non ho bisogno di un aggancio sincronizzato, ma ho bisogno di un filtro più stretto per evitare interferenze tra i canali adiacenti.
La tolleranza che potevo accettare era di 100 microsecondi, che è il tempo massimo che potevo tollerare per la sincronizzazione tra i terminali e la stazione base. Se la tolleranza fosse stata maggiore, avrei avuto problemi di interferenza tra i canali adiacenti, mentre se fosse stata minore, avrei avuto problemi di sincronizzazione e di perdita di pacchetti.\\

\paragraph{TDMA-FDMA}
Si usavano dei canali e il GSM era chiamato dual band perché potevo usare sia la banda a 900 MHz che quella a 1800 MHz. Ogni canale aveva una larghezza di banda di 200 kHz, e ogni canale era diviso in 8 timeslot, ognuno dei quali poteva essere assegnato a un utente diverso. In questo modo, ogni canale poteva supportare fino a 8 utenti contemporaneamente, e la capacità totale della rete era aumentata significativamente rispetto alla prima generazione.\\
\paragraph{Architettura di una rete 2G}
L'architettura di una rete 2G é rimasta. Il cellulare viene chiamato mobile terminal o mobile equipment, al suo interno ha la SIM che identifica l'utente. La connesione tra la Base station e la stazione base controller è chiamata Abis, mentre la connessione tra la stazione base controller e il mobile switching center è chiamata A. Il mobile switching center è responsabile della gestione delle chiamate e della connessione alla rete fissa, mentre la stazione base controller gestisce le risorse radio e la comunicazione con le stazioni base.\\
La Base station viene controllata da un BSC (Base Station Controller) che gestisce le risorse radio e la comunicazione con le stazioni base, mentre il MSC (Mobile Switching Center) è responsabile della gestione delle chiamate e della connessione alla rete fissa.\\
Devo avere informazioni relative all'home register e al visitor register, che sono database che contengono informazioni sugli utenti della rete. L'home register contiene informazioni sugli utenti registrati nella rete, mentre il visitor register contiene informazioni sugli utenti che si trovano temporaneamente nella rete, ad esempio quando si spostano da una cella all'altra. Queste informazioni sono utilizzate per gestire le chiamate e garantire la continuità del servizio quando gli utenti si spostano tra le celle.\\
EIR é l'informazione non sull'utente ma sul dispositivo che ha un database che contiene informazioni sui dispositivi mobili, come il numero di serie e il modello. Questo database viene utilizzato per identificare i dispositivi mobili e garantire la sicurezza della rete, ad esempio bloccando i dispositivi rubati o non autorizzati.\\

La SIm ha codice di sicurezza personali come PIN (Personal Identification Number) e PUK (Personal Unblocking Key) che sono utilizzati per proteggere l'accesso alla SIM e garantire la sicurezza delle comunicazioni. Il PIN è un codice a 4 cifre che l'utente deve inserire per accedere alla SIM, mentre il PUK è un codice a 8 cifre che viene utilizzato per sbloccare la SIM se il PIN viene inserito in modo errato più volte. Questi codici sono importanti per proteggere le informazioni dell'utente e garantire la sicurezza della rete.\\
IMEI é un codice univoco che identifica ogni dispositivo mobile, ed è utilizzato per garantire la sicurezza della rete e prevenire l'uso di dispositivi non autorizzati. Il codice IMEI viene assegnato a ogni dispositivo mobile durante la produzione e viene memorizzato nella SIM. Quando un dispositivo si connette alla rete, il codice IMEI viene verificato per garantire che il dispositivo sia autorizzato a utilizzare la rete. Se un dispositivo viene segnalato come rubato o non autorizzato, il codice IMEI può essere bloccato, impedendo al dispositivo di connettersi alla rete.\\

\paragraph{General Packet Radio Service - GPRS}
GPRS é un miglioramento del GSM che permette la trasmissione di dati a pacchetto, con velocità fino a 114 kbps. GPRS utilizza una tecnica di modulazione chiamata Gaussian Minimum Shift Keying (GMSK) e supporta sia la trasmissione dati che la trasmissione vocale. GPRS è stato un passo importante verso la 3G, poiché ha introdotto la possibilità di trasmettere dati a pacchetto, consentendo l'accesso a Internet e ad altri servizi basati sui dati. Tuttavia, GPRS aveva ancora limitazioni in termini di velocità di trasmissione dati, che sono state superate con l'introduzione di tecnologie come EDGE e UMTS.\\
Arriva da un internet service provider un indirizzo IP e il gateway deve inoltrare il pacchetto verso il base station. 

\subsection{3G}
La terza generazione di reti mobili, introdotta negli anni 2000, ha portato a un ulteriore miglioramento nella comunicazione mobile. Con l'introduzione del UMTS (Universal Mobile Telecommunications System), le velocità di trasmissione dati sono aumentate significativamente, consentendo l'accesso a Internet e ad altri servizi basati sui dati in modo più efficiente. Inoltre, il UMTS ha introdotto la possibilità di supportare servizi multimediali come videochiamate e streaming, migliorando ulteriormente l'esperienza dell'utente. Tuttavia, il UMTS aveva ancora limitazioni in termini di velocità di trasmissione dati, che sono state superate con l'introduzione di tecnologie come HSPA e LTE, che hanno portato alla transizione verso la 4G.\\
In italia il primo operatore a lanciare la rete 3G è stato 3 Italia, seguito da Vodafone e TIM. La copertura della rete 3G è stata inizialmente limitata alle aree urbane, ma si è rapidamente estesa a livello nazionale, consentendo a un numero sempre maggiore di utenti di accedere ai servizi basati sui dati. La rete 3G ha rappresentato un passo importante verso la connettività mobile avanzata, aprendo la strada a nuove applicazioni e servizi che hanno trasformato il modo in cui le persone comunicano e accedono alle informazioni.\\
\paragraph{Wideband Code Division Multiple Access - WCDMA}
WCDMA é una tecnica di accesso multiplo utilizzata nelle reti 3G, che consente a più utenti di condividere la stessa frequenza utilizzando codici univoci per distinguere i segnali. WCDMA utilizza una tecnica di modulazione chiamata Quadrature Phase Shift Keying (QPSK) e supporta sia la trasmissione dati che la trasmissione vocale. WCDMA è stato un passo importante verso la 4G, poiché ha introdotto la possibilità di trasmettere dati a pacchetto in modo più efficiente, consentendo l'accesso a Internet e ad altri servizi basati sui dati in modo più veloce e affidabile. Tuttavia, WCDMA aveva ancora limitazioni in termini di velocità di trasmissione dati, che sono state superate con l'introduzione di tecnologie come HSPA e LTE.\\
\paragraph{High Speed Packet Access - HSPA}
HSPA é un miglioramento del UMTS che permette velocità di trasmissione dati più elevate, fino a 14 Mbps. HSPA utilizza una tecnica di modulazione chiamata Quadrature Phase Shift Keying (QPSK) e supporta sia la trasmissione dati che la trasmissione vocale. HSPA è stato un passo importante verso la 4G, poiché ha introdotto la possibilità di trasmettere dati a pacchetto in modo più efficiente, consentendo l'accesso a Internet e ad altri servizi basati sui dati in modo più veloce e affidabile. Tuttavia, HSPA aveva ancora limitazioni in termini di velocità di trasmissione dati, che sono state superate con l'introduzione di tecnologie come LTE, che hanno portato alla transizione verso la 4G.\\
\paragraph{Architettura di una rete 3G}
L'architettura di una rete 3G è simile a quella di una rete 2G, ma con alcune modifiche per supportare le nuove tecnologie e i nuovi servizi. La rete 3G utilizza ancora stazioni base e stazioni base controller, ma introduce anche nuovi componenti come il Node B, che è la stazione base utilizzata nelle reti 3G, e il Radio Network Controller (RNC), che gestisce le risorse radio e la comunicazione tra i Node B. Inoltre, la rete 3G introduce un nuovo componente chiamato Serving GPRS Support Node (SGSN), che gestisce la trasmissione dei dati a pacchetto e la connessione alla rete fissa. L'architettura di una rete 3G è progettata per supportare servizi multimediali avanzati e velocità di trasmissione dati più elevate rispetto alle reti precedenti.\\
Lo User Equipment é il tewrminale dell'utente. Ora c'é la USIM (Universal Subscriber Identity Module) che è una versione avanzata della SIM, che supporta funzionalità aggiuntive come la crittografia e l'autenticazione avanzata. La USIM è utilizzata nelle reti 3G e successive per garantire la sicurezza delle comunicazioni e proteggere le informazioni dell'utente.\\
La RAN (Radio Access Network) è la parte della rete che gestisce la comunicazione tra i terminali mobili e le stazioni base. La RAN è composta da componenti come il Node B e il Radio Network Controller (RNC), che gestiscono le risorse radio e la comunicazione tra i terminali mobili e le stazioni base. La RAN è responsabile della gestione delle risorse radio, della trasmissione dei dati a pacchetto e della connessione alla rete fissa, garantendo una comunicazione efficiente e affidabile per gli utenti mobili.\\
IL nodo B converte il flusso dei dati digitali in segnali radio usando la modulazione WCDMA. \\
La parte di rete Core si divide in due parti, la parte di rete di commutazione (Circuit Switched - CS) che gestisce le chiamate vocali e la parte di rete a pacchetto (Packet Switched - PS) che gestisce la trasmissione dei dati a pacchetto. La parte di rete CS è responsabile della gestione delle chiamate vocali, mentre la parte di rete PS è responsabile della gestione della trasmissione dei dati a pacchetto e della connessione alla rete fissa.\\
I registri comuni sono rimasti piú o meno gli stessi, abbiamo ancora l'home register e il visitor register, ma abbiamo anche un nuovo registro chiamato Authentication Center (AuC), che è responsabile della gestione delle credenziali di autenticazione degli utenti e della generazione dei codici di autenticazione per garantire la sicurezza delle comunicazioni.\\

\paragraph{Multiple Input Multiple Output}
MIMO é una tecnologia utilizzata nelle reti 3G e successive che consente di utilizzare più antenne per trasmettere e ricevere segnali, migliorando la qualità del segnale e aumentando la capacità della rete. MIMO utilizza tecniche di codifica spaziale per trasmettere più flussi di dati contemporaneamente, consentendo di aumentare la velocità di trasmissione dati e migliorare l'affidabilità delle comunicazioni. MIMO è stato un passo importante verso la 4G, poiché ha introdotto la possibilità di trasmettere dati a pacchetto in modo più efficiente, consentendo l'accesso a Internet e ad altri servizi basati sui dati in modo più veloce e affidabile. Tuttavia, MIMO aveva ancora limitazioni in termini di velocità di trasmissione dati, che sono state superate con l'introduzione di tecnologie come LTE, che hanno portato alla transizione verso la 4G.\\
Tre tipi di MIMO: Spatioal Multiplexing, Diversity Coding e Beamforming. Il Spatial Multiplexing consente di trasmettere più flussi di dati contemporaneamente utilizzando più antenne, aumentando la velocità di trasmissione dati. Il Diversity Coding utilizza più antenne per trasmettere lo stesso flusso di dati, migliorando l'affidabilità delle comunicazioni in condizioni di segnale debole o interferenze. Il Beamforming utilizza più antenne per focalizzare il segnale verso un utente specifico, migliorando la qualità del segnale e aumentando la capacità della rete.\\

\subsection{Long Term Evolution}
La quarta generazione di reti mobili, introdotta negli anni 2010, ha portato a un ulteriore miglioramento nella comunicazione mobile. Con l'introduzione del LTE (Long Term Evolution), le velocità di trasmissione dati sono aumentate significativamente, consentendo l'accesso a Internet e ad altri servizi basati sui dati in modo più efficiente. Inoltre, il LTE ha introdotto la possibilità di supportare servizi multimediali avanzati come video in alta definizione e giochi online, migliorando ulteriormente l'esperienza dell'utente. Tuttavia, il LTE aveva ancora limitazioni in termini di velocità di trasmissione dati, che sono state superate con l'introduzione di tecnologie come 5G, che hanno portato alla transizione verso la 5G.\\
Le caratteristiche principali sono la mobilitá, sicurezza, roaming e prestazioni.
Per la mobilitá abbiamo la possibilità di spostarsi da una cella all'altra senza interruzioni, grazie al processo di handover che consente di trasferire la chiamata o la connessione dati da una stazione base a un'altra in modo fluido e senza interruzioni.\\
Per la sicurezza, il LTE utilizza meccanismi di autenticazione e crittografia avanzati per proteggere le comunicazioni e garantire la privacy degli utenti.\\
Per il roaming, il LTE supporta la possibilità di utilizzare i servizi mobili anche quando si è
fuori dalla propria rete domestica, consentendo agli utenti di connettersi a reti partner in tutto il mondo.\\
Per le prestazioni, il LTE offre velocità di trasmissione dati elevate, bassa latenza e una maggiore capacità della rete, consentendo agli utenti di accedere a servizi basati sui dati in modo più efficiente e affidabile.\\

MIMO é una tecnologia utilizzata nelle reti 4G e successive che consente di utilizzare più antenne per trasmettere e ricevere segnali, migliorando la qualità del segnale e aumentando la capacità della rete. MIMO utilizza tecniche di codifica spaziale per trasmettere più flussi di dati contemporaneamente, consentendo di aumentare la velocità di trasmissione dati e migliorare l'affidabilità delle comunicazioni. MIMO è stato un passo importante verso la 5G, poiché ha introdotto la possibilità di trasmettere dati a pacchetto in modo più efficiente, consentendo l'accesso a Internet e ad altri servizi basati sui dati in modo più veloce e affidabile. Tuttavia, MIMO aveva ancora limitazioni in termini di velocità di trasmissione dati, che sono state superate con l'introduzione di tecnologie come 5G, che hanno portato alla transizione verso la 5G.\\
SC-OFDMA é una tecnica di accesso multiplo utilizzata nelle reti 4G, che consente a più utenti di condividere la stessa frequenza utilizzando codici univoci per distinguere i segnali. SC-OFDM utilizza una tecnica di modulazione chiamata Orthogonal Frequency Division Multiplexing (OFDM) e supporta sia la trasmissione dati che la trasmissione vocale. SC-OFDM è stato un passo importante verso la 5G, poiché ha introdotto la possibilità di trasmettere dati a pacchetto in modo più efficiente, consentendo l'accesso a Internet e ad altri servizi basati sui dati in modo più veloce e affidabile. Tuttavia, SC-OFDM aveva ancora limitazioni in termini di velocità di trasmissione dati, che sono state superate con l'introduzione di tecnologie come 5G, che hanno portato alla transizione verso la 5G.\\
SDR (Software Defined Radio) é una tecnologia utilizzata nelle reti 4G e successive che consente di implementare le funzionalità radio in software, invece che in hardware. SDR consente di aggiornare e modificare le funzionalità radio in modo più flessibile e rapido, consentendo di adattarsi alle nuove tecnologie e ai nuovi standard senza dover sostituire l'hardware. SDR è stato un passo importante verso la 5G, poiché ha introdotto la possibilità di implementare nuove funzionalità radio in modo più efficiente, consentendo l'accesso a Internet e ad altri servizi basati sui dati in modo più veloce e affidabile. Tuttavia, SDR aveva ancora limitazioni in termini di velocità di trasmissione dati, che sono state superate con l'introduzione di tecnologie come 5G, che hanno portato alla transizione verso la 5G.\\
Scheduling Dinamico é una tecnica utilizzata nelle reti 4G e successive che consente di allocare le risorse radio in modo dinamico in base alle esigenze degli utenti e alle condizioni della rete. Lo scheduling dinamico utilizza algoritmi di scheduling per determinare quali utenti devono essere serviti in un dato momento, tenendo conto di fattori come la qualità del segnale, la quantità di dati da trasmettere e la priorità degli utenti. Lo scheduling dinamico è stato un passo importante verso la 5G, poiché ha introdotto la possibilità di allocare le risorse radio in modo più efficiente, consentendo l'accesso a Internet e ad altri servizi basati sui dati in modo più veloce e affidabile. Tuttavia, lo scheduling dinamico aveva ancora limitazioni in termini di velocità di trasmissione dati, che sono state superate con l'introduzione di tecnologie come 5G, che hanno portato alla transizione verso la 5G.\\
\paragraph{OFDM e OFDMA}
OFDM (Orthogonal Frequency Division Multiplexing) é una tecnica di modulazione utilizzata nelle reti 4G e successive che consente di trasmettere dati su più frequenze contemporaneamente, migliorando la velocità di trasmissione dati e l'efficienza spettrale. OFDM divide il canale di comunicazione in sottocanali ortogonali, consentendo di trasmettere più flussi di dati contemporaneamente senza interferenze. OFDMA (Orthogonal Frequency Division Multiple Access) é una tecnica di accesso multiplo basata su OFDM che consente a più utenti di condividere la stessa frequenza utilizzando codici univoci per distinguere i segnali. OFDMA è stato un passo importante verso la 5G, poiché ha introdotto la possibilità di trasmettere dati a pacchetto in modo più efficiente, consentendo l'accesso a Internet e ad altri servizi basati sui dati in modo più veloce e affidabile. Tuttavia, OFDMA aveva ancora limitazioni in termini di velocità di trasmissione dati, che sono state superate con l'introduzione di tecnologie come 5G, che hanno portato alla transizione verso la 5G.\\



\end{document}